\documentclass[11pt]{report}
\usepackage{fontspec}
\setmainfont{TeX Gyre Pagella}
\usepackage[margin=1.15in]{geometry}
\usepackage{amsmath,amssymb,amsthm}
\usepackage{enumitem}
\usepackage{array}
\usepackage{hyperref}
\hypersetup{colorlinks=true, linkcolor=black, urlcolor=black, citecolor=black}

\newtheorem{proposition}{Proposition}[chapter]
\newtheorem{theorem}{Theorem}[chapter]
\newtheorem{lemma}{Lemma}[chapter]
\newtheorem{definition}{Definition}[chapter]
\newtheorem{remark}{Remark}[chapter]
\newtheorem{corollary}{Corollary}[chapter]
\newtheorem{example}{Example}[chapter]
\newtheorem{openproblem}{Open Problem}[chapter]
\newtheorem{caveat}{Caveat}[chapter]

\title{\textbf{Continuation Geometry}}
\author{Flyxion}
\date{July 2026}

\begin{document}

\maketitle

\tableofcontents

\chapter{Why Continuation Matters}

A system that has stopped changing has not necessarily stopped existing,
but a system that has run out of ways to change has. There is a
difference between a room that is quiet because nothing needs to happen
in it and a room that is quiet because nothing can. The first is rest.
The second is the end of a story that may not announce itself as an
ending. Most of what concerns us about a living body, a piece of
software, an institution, or a plan is not its condition at the present
instant but the width of what it can still become. A body in a coma and
a body in deep sleep can look alike from across the room; what
distinguishes them is not visible in a single frame but in the shape of
what still lies ahead of each.

This is a book about that shape. What a system could still do matters
more than what it is doing right now: the whole reachable fan of
futures opening out from wherever it happens to stand. Call that fan a
system's futures, or its options, or its story-space. Whatever you
call it, it's the thing that shrinks when a hinge rusts shut, when a
species loses a gene it won't get back, when an argument commits
itself to a position it can't walk back from, when a bridge corrodes
past the point where any repair schedule could save it. Repair, in the
most general sense of the word, tries to restore it: not a particular
outcome, but the standing
possibility of many outcomes.

Several habits of thought already study pieces of this. Engineers ask
whether a structure will hold under load. Ecologists ask whether a
population can recover from a disturbance. Economists ask whether an
allocation can be sustained, and control theorists ask whether a
system can be steered to stay within safe bounds indefinitely. Each of
these is really the same question in different clothes: whether the
fan of futures remains nonempty, whether there is still at least one
way forward that respects whatever constraints the situation imposes.
That question has a mature and exact mathematical answer, developed
over several decades under the name of viability theory, and this book
leans on that answer throughout.

But existence is not the only thing worth asking about a fan of futures.
Two systems can both have somewhere left to go and still differ enormously
in how much room they have to maneuver, how easily they could be nudged
back from the edge, which of their internal distinctions survive into
their futures and which quietly get erased, and which of their many
possible continuations is the one actually worth choosing. A patient who
can still recover and a patient who can barely still recover are both,
technically, viable; no doctor would treat that technicality as the end
of the conversation. The questions that begin where mere existence leaves off, how much
room, how easily recovered, what is preserved, what is worth choosing,
are what this book is organized around. They
do not yet have as settled a mathematical treatment as the existence
question does. Building one, on top of the existence question rather than
instead of it, is the project.

The order of concerns runs roughly as follows, and will recur throughout
the book. A system draws distinctions: it treats some states as
meaningfully different from others. Those distinctions are carried by
information, and information degrades. This is what entropy names.
Some of that degradation can be undone; the undoing is repair. Repair is
only possible where a future still exists to be repaired into; that
existence is continuation. And among the continuations that exist, some
are better than others by standards that have nothing to do with mere
existence; sorting them out is admissibility. Existence, in other words,
is necessary for everything downstream of it but is far from sufficient
for any of it. This book starts at existence, because it is the
best-understood link in the chain, and spends the rest of its length
building outward from it toward repair, distinction, and admissibility:
the concerns that motivated picking up the chain in the first place.

\section*{The object of study}

We now fix, once and for all, the formal object the rest of the book is
about.

\begin{definition}[State, dynamics, constraint]
Let $X$ be a state space (a metric space; in applications typically
$\mathbb{R}^n$ or a finite-dimensional manifold). Let
$F : X \rightrightarrows X$ be a set-valued map assigning to each state
$x \in X$ a set $F(x) \subseteq X$ of admissible velocities, so that
trajectories of the system are absolutely continuous curves
$\gamma$ satisfying the differential inclusion
$\dot\gamma(t) \in F(\gamma(t))$. Let $K \subseteq X$ be closed; $K$ is
the constraint set, the region within which the system is required to
remain.
\end{definition}

\begin{definition}[Continuation space]
For $x \in K$, the continuation space of $x$ is
\[
\Gamma_K(x) = \{\gamma : [0,\infty) \to K \mid \gamma(0) = x,\ \dot\gamma(t) \in F(\gamma(t))\ \text{for a.e.\ } t \geq 0\}.
\]
For a finite horizon $T > 0$, the finite-horizon continuation space is
\[
\Gamma_K^T(x) = \{\gamma : [0,T] \to K \mid \gamma(0) = x,\ \dot\gamma(t) \in F(\gamma(t))\ \text{for a.e.\ } t \in [0,T]\}.
\]
\end{definition}

$\Gamma_K(x)$ (or simply $\Gamma(x)$ where $K$ is fixed by context) is
the fan of futures referred to above, made precise relative to one
particular reading of ``possible'': a future is possible if it is a
trajectory of the given dynamics that never leaves the given constraint
set. Other readings of possible are available and will matter later:
controllably possible, admissibly possible, epistemically possible. But
this dynamical reading is where the book begins, because it is the one
with an existing, exact theory attached to it.

\begin{proposition}[Restriction is well-defined]
For every $T > 0$, restriction of a trajectory to $[0,T]$ defines a map
\[
\rho_T : \Gamma_K(x) \to \Gamma_K^T(x), \qquad \rho_T(\gamma) = \gamma|_{[0,T]}.
\]
\end{proposition}
\begin{proof}
If $\gamma \in \Gamma_K(x)$ then $\gamma(0) = x$, $\gamma(t) \in K$ for
all $t \geq 0$, and $\dot\gamma(t) \in F(\gamma(t))$ for a.e.\ $t \geq 0$.
Restricting to $[0,T]$ preserves each of these three conditions on the
smaller domain, so $\gamma|_{[0,T]} \in \Gamma_K^T(x)$.
\end{proof}

\begin{proposition}[Consistency across horizons]
For $0 < T_1 < T_2$, restriction further to $[0,T_1]$ defines a map
$\rho_{T_1, T_2} : \Gamma_K^{T_2}(x) \to \Gamma_K^{T_1}(x)$, and the
diagram of restriction maps commutes:
$\rho_{T_1} = \rho_{T_1,T_2} \circ \rho_{T_2}$ for every
$T_1 < T_2 < \infty$.
\end{proposition}
\begin{proof}
Immediate from the definition of restriction: restricting a curve to
$[0,T_1]$ gives the same result whether it is done directly from the
domain $[0,\infty)$ or via an intermediate restriction to $[0,T_2]$,
since restriction of a function to a subset of its domain depends only on
that subset, not on the route by which the subset was reached.
\end{proof}

\begin{remark}
This consistency is what allows $\Gamma_K(x)$ to be recovered, in good
cases, as the projective (inverse) limit of the finite-horizon spaces
$\Gamma_K^T(x)$ as $T \to \infty$. This observation is recorded here
because Chapter 6 will build continuation volume and continuation
distance on the finite-horizon spaces precisely because they are
well-posed where the infinite-horizon trajectory space is not; the
present proposition is what justifies treating the finite-horizon
construction as an approximation to $\Gamma_K(x)$ rather than a
different object entirely.
\end{remark}

\begin{proposition}[Equilibria have nonempty continuation spaces]
If $x^\ast \in K$ satisfies $0 \in F(x^\ast)$, then $\Gamma_K(x^\ast) \neq \varnothing$.
\end{proposition}
\begin{proof}
Let $\gamma(t) = x^\ast$ for all $t \geq 0$. Then $\gamma(0) = x^\ast$,
$\gamma(t) = x^\ast \in K$ for all $t$, and $\dot\gamma(t) = 0 \in F(x^\ast) = F(\gamma(t))$
for all $t$. Hence $\gamma \in \Gamma_K(x^\ast)$.
\end{proof}

\begin{remark}
This is the first, and simplest, example of a nonempty continuation
space, and it is worth pausing on because it already illustrates a
distinction the rest of the book will not let go of: nonemptiness of
$\Gamma_K(x^\ast)$ here is witnessed by a single, unmoving trajectory. It
says nothing about how many other continuations $x^\ast$ has, how close
$x^\ast$ sits to states with no continuations at all, or whether the
distinctions the system cares about survive along the witnessing
trajectory. A full characterization of nonemptiness for arbitrary
$x \in K$, not just equilibria, is the content of Aubin's Viability
Theorem, taken up next.
\end{remark}

\section*{The progression this book follows}

The chain of concerns described in natural language above, distinction,
information, entropy, repair, continuation, admissibility, is not a
theorem and is not offered as one. It is a roadmap, and its purpose is to
say in advance which part of the whole problem each part of the book is
responsible for.

Continuation, in the sense of $\Gamma_K(x) \neq \varnothing$, is the
best-understood link and is addressed first: it is the subject of
viability theory, surveyed and put to use in the chapters immediately
following this one. Repair is addressed once continuation is in hand,
since a repair is, formally, nothing more than a continuation that
reaches a designated target: the machinery for saying this precisely
is capture, and it costs nothing beyond what continuation has already
supplied. Distinction and admissibility are addressed last, and are the
hardest part of the roadmap to formalize, because they require deciding
not merely whether a future exists but which of the many existing
futures should be preferred, and by what standard. The book's later
chapters build that standard in stages rather than assuming it in
advance.

\chapter{Precursors}

No idea arrives without ancestors, and an idea that pretends otherwise
usually turns out to be a poorer version of something already on the
shelf. The temptation for any project that finally finds its central
object is to treat the discovery as an arrival rather than a
continuation of its own, to write the history chapter last, briefly,
and mostly as a courtesy. That temptation is worth resisting here for a
specific reason: the people who spent real years on the questions this
book asks were not working on one shared problem that only now gets
its proper name. They were working on at least two different problems
that happen, from a certain distance, to look like the same problem.

The first tradition asks how a system manages to keep going. Given
where a system stands and how it is allowed to move, will it still be
able to move tomorrow, and the day after, without breaking whatever
rule keeps it a functioning version of itself? This is a question about
persistence, and it has been asked with increasing precision for over a
century: orbits that never quite settle down, equilibria that recover
from being nudged, decisions that must be made in sequence without
knowing the whole future in advance. Each advance in this tradition
sharpened the same underlying question rather than replacing it, which
is why it reads, in retrospect, like a single conversation carried on
by different people.

The second tradition asks a different question, one that does not
reduce to the first even though it sounds similar. Given where a system
stands, what does the whole space of things it could still become look
like? Not whether that space is empty, but what shape it has, how it
grows or shrinks as choices get made, and which of its features are
worth preserving. This tradition is younger, more scattered across
disciplines that rarely cite one another, and considerably less tidy.
Evolving biological form. What innovations become available once other
innovations have already happened. Causal alternatives to what
actually occurred. Even an architect's attempt to say which building
layouts leave a neighborhood room to keep being a good place to live.
None of these projects thought of itself as contributing to a general
theory of possibility-spaces. Read together, they look like exactly
that.

This book sits where these two traditions meet, and owes each of them a
debt it should not obscure. What follows makes some of that debt
precise, showing in each case how much of the earlier work can be
recovered as a special case or a direct ancestor of the object this
book is built around, and where each earlier tradition stops short.

\section*{The persistence lineage}

\begin{definition}[Recurrent state, after Poincar\'e]
Let $\phi_t : X \to X$ be a measure-preserving flow on a probability
space $(X,\mu)$, and let $A \subseteq X$ have $\mu(A) > 0$. A point
$x \in A$ is recurrent in $A$ if $\phi_t(x) \in A$ for some $t > 0$,
and infinitely often for a.e.\ such $x$.
\end{definition}

\begin{proposition}[Poincar\'e recurrence and nonempty continuation]
If $A$ is invariant under $\phi_t$ in the sense that $\phi_t^{-1}(A)$
has full relative measure in $A$ for all $t$, then for a.e.\ $x \in A$,
the trajectory $\gamma(t) = \phi_t(x)$ lies in $\Gamma_A(x)$, and
moreover returns to $A$ infinitely often.
\end{proposition}
\begin{proof}
This is the Poincar\'e Recurrence Theorem, applied to the flow
restricted to $A$: almost every point of a measure-preserving invariant
set returns to that set infinitely often under forward iteration
\cite{poincare1890}. The trajectory $\gamma(t) = \phi_t(x)$ satisfies
$\gamma(0) = x$ and $\gamma(t) \in A$ for all $t \geq 0$ by invariance
of $A$, and $\dot\gamma(t)$ equals the generator of $\phi_t$ at
$\gamma(t)$, which lies in $F(\gamma(t))$ trivially when $F$ is taken to
be the vector field generating $\phi_t$. Hence $\gamma \in \Gamma_A(x)$.
\end{proof}

\begin{remark}
This is the earliest instance, in the lineage this book claims, of a
proof that a continuation space is nonempty, though Poincar\'e's
concern was not persistence under constraint but the much stronger fact
of eventual return. Recurrence is strictly stronger than mere viability:
every recurrent trajectory witnesses nonemptiness of $\Gamma_A(x)$, but
nonemptiness of $\Gamma_K(x)$ does not require any trajectory to return
to its starting neighborhood at all. The relation between recurrence and
viability, when a viable trajectory can be strengthened to a recurrent
one, is not pursued in this book but is flagged here as a legitimate
question left on the table by folding Poincar\'e into the same lineage
as Aubin.
\end{remark}

\begin{definition}[Lyapunov stability]
A state $x^\ast$ with $0 \in F(x^\ast)$ is Lyapunov stable if for every
$\varepsilon > 0$ there exists $\eta > 0$ such that
$\|x - x^\ast\| < \eta$ implies every $\gamma \in \Gamma_X(x)$ satisfies
$\|\gamma(t) - x^\ast\| < \varepsilon$ for all $t \geq 0$.
\end{definition}

\begin{proposition}[Lyapunov functions certify viability of sublevel sets]
Suppose $V : X \to \mathbb{R}_{\geq 0}$ is continuously differentiable,
$V(x^\ast) = 0$, $V(x) > 0$ for $x \neq x^\ast$, and
$\sup_{v \in F(x)} \nabla V(x) \cdot v \leq 0$ for all $x$. Then for
every $c > 0$, the sublevel set $K_c = \{x : V(x) \leq c\}$ is viable
under $F$: $\Gamma_{K_c}(x) \neq \varnothing$ for every $x \in K_c$.
\end{proposition}
\begin{proof}
Fix $x \in K_c$ and let $\gamma$ solve $\dot\gamma(t) \in F(\gamma(t))$,
$\gamma(0) = x$ (such a solution exists locally under standard
regularity hypotheses on $F$). Along $\gamma$,
$\frac{d}{dt} V(\gamma(t)) = \nabla V(\gamma(t)) \cdot \dot\gamma(t) \leq
\sup_{v \in F(\gamma(t))} \nabla V(\gamma(t)) \cdot v \leq 0$, so
$V(\gamma(t))$ is nonincreasing. Since $V(\gamma(0)) = V(x) \leq c$, it
follows that $V(\gamma(t)) \leq c$ for all $t$, i.e.\ $\gamma(t) \in K_c$
for all $t$. Hence $\gamma \in \Gamma_{K_c}(x)$.
\end{proof}

\begin{remark}
This is the classical Lyapunov stability argument \cite{lyapunov1892},
restated to make explicit what it already was: a direct construction of
a nonempty, in fact tangentially controlled, continuation space. The
tangency criterion of Aubin's Viability Theorem can be read as the
generalization of this proposition obtained by dropping the requirement
that a single global function $V$ certify viability for every sublevel
set at once, and instead checking the weaker, purely local condition
$F(x) \cap T_K(x) \neq \varnothing$ directly on the boundary of whatever
set $K$ is given, without needing $K$ to arise as a sublevel set of
anything.
\end{remark}

\begin{definition}[Bellman's principle of optimality]
For a controlled system with running cost $l$ and terminal cost $c$,
define the value function
\[
J(t,x) = \inf_{u(\cdot)} \left[ \int_t^T l(x(s), u(s))\, ds + c(x(T)) \right]
\]
over admissible controls $u(\cdot)$ steering $\dot x(s) = f(x(s), u(s))$
from $x(t) = x$.
\end{definition}

\begin{proposition}[Dynamic programming as a special case of guaranteed valuation]
Suppose $Q(x) = \{0\}$ (no perturbation) in Aubin's dynamic core
construction. Then Aubin's guaranteed valuation function $V^\sharp(T,x)$
reduces to Bellman's value function $J(t,x)$ under the correspondence
$T \leftrightarrow T - t$, and the Hamilton-Jacobi-Isaacs variational
inequality reduces to the Hamilton-Jacobi-Bellman equation
\[
-\frac{\partial J}{\partial t} + \inf_{u \in P(x)}\left[ \nabla J(t,x) \cdot f(x,u) + l(x,u) \right] = 0.
\]
\end{proposition}
\begin{proof}
When $Q(x) = \{0\}$, the supremum over perturbations $v \in Q(x)$ in the
defining formula for $V^\sharp$ (Aubin, Theorem 1.2) has a single term
to range over, so it drops out, leaving only the infimum over feedbacks
$p_e(x) \in P(x)$. This is exactly Bellman's infimum over controls, and
the discount term $m(x,p,v)$, in the absence of perturbation, plays the
role of a running cost once $l$ is identified with the appropriate
component of the functional $J_u$. Substituting into the
Hamilton-Jacobi-Isaacs inequality and dropping the now-vacuous
$\sup_{v \in Q(x)}$ term gives the stated Hamilton-Jacobi-Bellman
equation \cite{bellman1957}.
\end{proof}

\begin{remark}
Dynamic programming is, on this reading, the special case of guaranteed
valuation in which nature is not permitted to intervene. Aubin's
contribution to the persistence lineage was to reinstate the adversary
that Bellman's formulation sets to zero, without discarding the
optimization structure Bellman had already built. Aubin's viability
kernel and capture basin are the further generalization obtained by
dropping the requirement that a single scalar cost be optimized at all,
and asking only whether \emph{some} admissible control keeps the system
where it needs to be: existence in place of optimality, which is
where this book's own starting point sits.
\end{remark}

\section*{The possibility-space lineage}

\begin{definition}[Adjacent possible, after Kauffman]
Given a set of currently realized configurations $S_t \subseteq X$ and a
one-step generative map $\Phi : X \to 2^X$ (the set of configurations
directly reachable from a given one), the adjacent possible of $S_t$ is
\[
\mathrm{Adj}(S_t) = \bigcup_{x \in S_t} \Phi(x).
\]
\end{definition}

\begin{proposition}[Adjacent-possible growth as reachable-set recursion]
Define $S_{t+1} = S_t \cup \mathrm{Adj}(S_t)$. Then, taking
$\Phi$ as a discrete-time one-step reachability map and $t$ discrete
steps,
\[
S_t = \bigcup_{x \in S_0} R_t(x),
\]
where $R_t(x)$ is the discrete-time finite-horizon reachable set
generated by $\Phi$ from $x$.
\end{proposition}
\begin{proof}
By induction on $t$. The base case $t=0$ is immediate:
$R_0(x) = \{x\}$ and $S_0 = \bigcup_{x \in S_0} \{x\}$. For the
inductive step, assume $S_t = \bigcup_{x \in S_0} R_t(x)$. Then
\[
S_{t+1} = S_t \cup \mathrm{Adj}(S_t) = \bigcup_{x \in S_0} R_t(x) \cup \bigcup_{y \in S_t} \Phi(y) = \bigcup_{x \in S_0} \left( R_t(x) \cup \bigcup_{y \in R_t(x)} \Phi(y) \right) = \bigcup_{x \in S_0} R_{t+1}(x),
\]
using that $\Phi(y)$ for $y \in R_t(x)$ is exactly the one-step
extension of a $t$-step reachable trajectory from $x$, which is what
$R_{t+1}(x)$ collects.
\end{proof}

\begin{remark}
Kauffman's adjacent possible \cite{kauffman2000} is, under this
correspondence, an early and purely combinatorial instance of the
finite-horizon reachable-set construction that later chapters use to
build continuation volume. What it lacks, by design, is any constraint
set $K$: the adjacent possible is not asked to remain within bounds,
only to grow, so it corresponds to the unconstrained case $K = X$. It is
this unconstrained framing, uncommon in the persistence lineage but
common in Kauffman's and Holland's own writing \cite{holland1995}, that
marks the possibility-space lineage as a genuinely separate tradition
rather than a duplicate of the first.
\end{remark}

\begin{definition}[Structural causal model and intervention, after Pearl]
A structural causal model is a tuple $(V, U, \mathcal{F})$ of endogenous
variables $V$, exogenous variables $U$, and structural equations
$\mathcal{F}$, each $f_i \in \mathcal{F}$ determining $V_i$ as a function
of its parents and relevant exogenous noise. An intervention
$\mathrm{do}(V_j = v)$ replaces $f_j$ with the constant function $v$,
producing a modified model $(V, U, \mathcal{F}_{\mathrm{do}(V_j=v)})$.
\end{definition}

\begin{proposition}[Counterfactual reachability as continuation under intervened dynamics]
Let $F$ correspond to the structural equations $\mathcal{F}$ in the
sense that $F(x)$ encodes the admissible next-states consistent with
$\mathcal{F}$ from configuration $x$, and let $F_{\mathrm{do}(V_j=v)}$
be the corresponding map under the intervened model. Then the
counterfactual ``what would have happened to $x$ had $V_j$ been set to
$v$'' corresponds exactly to $\Gamma_K^{F_{\mathrm{do}(V_j=v)}}(x)$, the
continuation space of $x$ computed under the intervened dynamics rather
than the original one.
\end{proposition}
\begin{proof}
By construction: a counterfactual trajectory in Pearl's sense is
obtained by holding the exogenous noise $U$ fixed at its actual values
and re-evaluating the structural equations under the modified
$\mathcal{F}_{\mathrm{do}(V_j=v)}$ \cite{pearl2000}. This is precisely
the definition of a trajectory of the continuation space
$\Gamma_K^{F_{\mathrm{do}(V_j=v)}}(x)$: an admissible evolution from $x$
under the dynamics $F_{\mathrm{do}(V_j=v)}$ rather than $F$.
\end{proof}

\begin{remark}
This is the cleanest correspondence in the possibility-space lineage,
and it points at something later chapters will need directly: comparing
continuations across different choices of dynamics $F$, not merely
across different starting states $x$ under one fixed $F$. Every
construction introduced so far has been stated for a single fixed
$(F,K)$. Pearl's framework is a reminder that the more general
comparison, across interventions on $F$ itself, is both meaningful and
already partly worked out elsewhere; this book does not take it up in
full, but the treatment of multiple agents sharing a dynamics later on
is the nearest point of contact.
\end{remark}

\begin{remark}[Levin: morphological and bioelectric continuation]
Levin's work on bioelectric signaling as a substrate for the storage and
recall of anatomical target-morphologies \cite{levin2021} is included in
this lineage on different grounds than the constructions above: not
because it yields a theorem restatable in the $\Gamma_K(x)$ vocabulary
without loss, but because it is the clearest existing example of a
system whose continuation space genuinely depends on more than its
instantaneous state: a regenerating organism's available futures
depend on bioelectric pattern memory that persists through and beyond
local tissue damage, in a way not recoverable from cell-by-cell
instantaneous configuration alone. This is, in the vocabulary already
fixed, exactly the history-dependence flagged as an open problem for the
continuation-operator construction. No proposition is offered here
forcing Levin's biology into the present formalism; the honest claim is
narrower, that this is the strongest existing example motivating why
history-dependence needs to be taken seriously rather than assumed away.
\end{remark}

\begin{remark}[Alexander: pattern languages]
Christopher Alexander's pattern language \cite{alexander1977} is the
loosest member of this lineage and is kept loose deliberately. A
pattern, in Alexander's sense, is a reusable solution to a recurring
design problem, stated together with the conditions under which it
applies and the further patterns it tends to invite or foreclose. Read
architecturally rather than dynamically, a choice of pattern is a move
that changes which further patterns remain available, which is a
possibility-space claim in substance, even though Alexander never
formalized it as one and this book will not retroactively supply the
formalization on his behalf. What is worth taking from Alexander is not
a theorem but a discipline: that a possibility-space can be studied
through the patterns that compose it and the compatibility relations
between them, prior to and independent of any measure, metric, or
dynamics being placed on it. That discipline anticipates the
lattice-based treatment of admissibility taken up later, where
composability and compatibility, rather than a single scalar score, are
again the primary structure.
\end{remark}

\section*{Where the lineages meet}

The two lineages traced here answer different questions and were not, in
general, in conversation with one another. The persistence lineage
answers whether a future exists; its instruments, recurrence,
stability, dynamic programming, viability, get sharper across a
century of work without changing what they are fundamentally asking. The
possibility-space lineage answers what the space of futures looks like
once one exists; its instruments are more heterogeneous because the
question is less settled, and in at least two cases above (Levin,
Alexander) the correspondence to the present formalism is offered as
motivation rather than as an established equivalence. The chapters that
follow begin by finishing what the persistence lineage started: stating
Aubin's Viability Theorem properly and using it as the foundation
everything else stands on, before turning to the harder task of giving
the possibility-space lineage's questions the same kind of precision
the persistence lineage has already earned.

\chapter{Viability Theory}

Walk along the top of a sea wall and there is a moment, well before you
lose your footing, when the only directions left that keep you on the
wall are a narrower and narrower band around straight ahead. You do not
fall the instant that band closes to nothing; you fall the instant
after, when a gust or a stumble sends you in some direction the band no
longer contains. Whether you are in danger, in other words, is not
really a question about your current position. Position alone cannot
distinguish the middle of a wide, flat wall from its crumbling edge if
both happen to put your feet at the same height. What matters is
narrower than position: it is whether the directions physics is about
to hand you and the directions that keep you on the wall still overlap.

This is a more useful fact than it first appears, because it turns a
question about the entire rest of your walk (will I, at any point in
the future, stumble off this wall) into a question you can answer by
looking only at where your feet are right now and which way you could
step next. The whole of the future collapses into a local, immediate
test. That collapse is not a metaphor; it is a theorem, discovered
first for smooth boundaries by Nagumo in the 1940s and given its
general, set-valued form by Aubin and Frankowska decades later. It is
the single fact this book leans on hardest, because everything from
here on, volume, distance, repair, admissibility, is built on top of a
boundary between what remains possible and what does not, and this is
the theorem that tells us how to find that boundary without having to
simulate every possible future first.

\section*{Standing hypotheses}

Throughout this chapter, $X = \mathbb{R}^n$, $K \subseteq X$ is closed,
and $F : X \rightrightarrows X$ is a set-valued map satisfying the
following standing hypotheses, in force unless stated otherwise:

\begin{enumerate}[label=(H\arabic*)]
\item \emph{Nonempty compact convex values.} $F(x)$ is nonempty,
compact, and convex for every $x \in X$.
\item \emph{Upper semicontinuity.} For every $x \in X$ and every open
$U \supseteq F(x)$, there is a neighborhood $V$ of $x$ such that
$F(y) \subseteq U$ for all $y \in V$.
\item \emph{Linear growth.} There is $c > 0$ such that
$\sup\{\|v\| : v \in F(x)\} \leq c(1 + \|x\|)$ for all $x \in X$.
\end{enumerate}

\begin{remark}
These are the hypotheses under which the classical existence theory for
differential inclusions, and hence the theorem below, holds without
further qualification \cite{aubincellina1984}. They are restrictive
enough to rule out pathological $F$, and permissive enough to cover
every example used in this book.
\end{remark}

\section*{Contingent cones}

\begin{definition}[Contingent cone]
For $K \subseteq X$ closed and $x \in K$, the contingent (Bouligand)
cone to $K$ at $x$ is
\[
T_K(x) = \left\{ v \in X : \liminf_{h \to 0^+} \frac{\operatorname{dist}(x + hv, K)}{h} = 0 \right\}.
\]
Equivalently, $v \in T_K(x)$ iff there exist sequences $h_n \to 0^+$ and
$v_n \to v$ with $x + h_n v_n \in K$ for all $n$.
\end{definition}

\begin{proposition}[Contingent cone at a smooth boundary]
Suppose $K = \{x : g(x) \leq 0\}$ for some continuously differentiable
$g : X \to \mathbb{R}$ with $\nabla g(x) \neq 0$ whenever $g(x) = 0$.
Then for $x$ with $g(x) = 0$,
\[
T_K(x) = \{v \in X : \nabla g(x) \cdot v \leq 0\}.
\]
\end{proposition}
\begin{proof}
($\supseteq$) Suppose $\nabla g(x) \cdot v \leq 0$. By Taylor expansion,
$g(x + hv) = g(x) + h \nabla g(x) \cdot v + o(h) = h \nabla g(x) \cdot v + o(h) \leq o(h)$,
so for $h$ small enough $g(x+hv) \leq 0$ up to an $o(h)$ correction,
giving $\operatorname{dist}(x+hv, K) = o(h)$ and hence $v \in T_K(x)$.

($\subseteq$) Suppose $\nabla g(x) \cdot v > 0$. Then by the same Taylor
expansion, $g(x+hv) = h\nabla g(x)\cdot v + o(h) > 0$ for $h$ small
enough, and since $\nabla g(x) \neq 0$, the distance from $x+hv$ to
$\{g \leq 0\}$ is bounded below by a constant multiple of
$g(x+hv)$ for $h$ small, which is bounded below by a constant multiple
of $h$. Hence $\operatorname{dist}(x+hv,K)/h$ does not tend to $0$, so
$v \notin T_K(x)$.
\end{proof}

\begin{remark}
This recovers the Lyapunov-style computation of the previous chapter as
the special case $g = V - c$ on the sublevel set $K_c = \{V \leq c\}$:
the tangency condition $F(x) \cap T_K(x) \neq \varnothing$ becomes
$\min_{v \in F(x)} \nabla V(x) \cdot v \leq 0$, exactly the condition
used there. The Viability Theorem below is what allows this
computation to be done boundary point by boundary point, without
requiring a single function $V$ to work for the whole of $K$ at once.
\end{remark}

\section*{The Viability Theorem}

\begin{theorem}[Aubin's Viability Theorem]
Under the standing hypotheses, $K$ is viable under $F$, that is,
$\operatorname{Viab}_F(K) = K$, if and only if
\[
F(x) \cap T_K(x) \neq \varnothing \qquad \text{for every } x \in K.
\]
\end{theorem}

\begin{proof}[Proof of necessity]
Suppose $\operatorname{Viab}_F(K) = K$, and fix $x \in K$. Then
$\Gamma_K(x) \neq \varnothing$; let $\gamma \in \Gamma_K(x)$. For a.e.\
$t$, $\dot\gamma(t) \in F(\gamma(t))$; since $\gamma$ is absolutely
continuous, it is differentiable at a.e.\ point, and in particular there
is a sequence $h_n \to 0^+$ of times at which $\gamma$ is differentiable
with $\dot\gamma(0^+) := \lim_{h_n \to 0^+} \frac{\gamma(h_n) - x}{h_n}$
existing and lying in $F(x)$ by upper semicontinuity of $F$ and
continuity of $\gamma$ at $0$. Since $\gamma(h_n) \in K$ for every $n$
(as $\gamma \in \Gamma_K(x)$), the sequence $v_n = \frac{\gamma(h_n)-x}{h_n}$
satisfies $x + h_n v_n = \gamma(h_n) \in K$, which is exactly the
sequential definition of membership in $T_K(x)$ for the limit
$v = \dot\gamma(0^+)$. Hence $v \in F(x) \cap T_K(x)$, so the
intersection is nonempty.
\end{proof}

\begin{proof}[Proof sketch of sufficiency]
Suppose $F(x) \cap T_K(x) \neq \varnothing$ for every $x \in K$; fix
$x_0 \in K$. We construct a viable trajectory from $x_0$ by an Euler
polygon argument, following the classical method \cite{aubin1991}.

Fix $\varepsilon > 0$ and a step size $h > 0$. At $x_0$, choose
$v_0 \in F(x_0) \cap T_K(x_0)$; by definition of the contingent cone,
there is a point within distance $\varepsilon h$ of $x_0 + h v_0$
lying in $K$; call it $x_1$. Repeat from $x_1$: choose
$v_1 \in F(x_1) \cap T_K(x_1)$, find $x_2 \in K$ within $\varepsilon h$
of $x_1 + h v_1$, and so on, producing a sequence $x_0, x_1, x_2, \dots
\in K$ and a piecewise-linear curve $\gamma_{h,\varepsilon}$ interpolating
them.

Under the linear growth hypothesis (H3), a discrete Gr\"onwall argument
bounds $\|x_k\|$ uniformly on any fixed time interval $[0,M]$,
independently of $h$ and $\varepsilon$ (for $h,\varepsilon$ small
enough), which in turn bounds $\|v_k\|$ uniformly on $[0,M]$ by (H3)
again. The curves $\gamma_{h,\varepsilon}$ are therefore uniformly
bounded and equi-Lipschitz on $[0,M]$ for each fixed $M$, so by the
Arzel\`a--Ascoli theorem, a subsequence converges uniformly on $[0,M]$
as $h, \varepsilon \to 0$; a diagonal argument over $M = 1,2,3,\dots$
produces a single limit curve $\gamma : [0,\infty) \to X$.

Since $K$ is closed and each $x_k \in K$ with interpolation error
vanishing as $\varepsilon \to 0$, the limit satisfies $\gamma(t) \in K$
for all $t \geq 0$. Since $F$ has closed graph (a consequence of upper
semicontinuity together with compact values) and each $F(x)$ is convex,
a standard weak-convergence argument (Mazur's lemma, applied to the
difference quotients of $\gamma_{h,\varepsilon}$) upgrades the a.e.\
subsequential convergence of the derivatives to
$\dot\gamma(t) \in F(\gamma(t))$ for a.e.\ $t$. Hence
$\gamma \in \Gamma_K(x_0)$, so $x_0 \in \operatorname{Viab}_F(K)$.

The technical steps compressed here, the discrete Gr\"onwall bound and
the passage from weak convergence of difference quotients to inclusion
in $F$ via Mazur's lemma, are standard in the differential inclusions
literature and are given in full in \cite{aubincellina1984} and
\cite{aubinframkowska1990}.
\end{proof}

\begin{remark}
Necessity did not need convexity of $F(x)$; sufficiency does, precisely
at the step converting weak convergence of difference quotients into
membership in $F(\gamma(t))$. This asymmetry is typical of existence
theorems for differential inclusions and is one reason (H1) is stated
as a standing hypothesis rather than derived.
\end{remark}

\section*{Finite horizons and the viability kernel}

\begin{definition}
$V_T(K) = \{x \in K : \Gamma_K^T(x) \neq \varnothing\}$, the set of
states admitting a trajectory remaining in $K$ for at least time $T$.
\end{definition}

By Proposition~1 of the opening chapter (restriction is well-defined),
$V_{T_2}(K) \subseteq V_{T_1}(K)$ whenever $T_1 \leq T_2$: a trajectory
valid on a longer horizon restricts to one valid on a shorter horizon,
so the finite-horizon sets are nested and shrinking in $T$.

\begin{theorem}[Viability kernel as a nested intersection]
Under the standing hypotheses,
\[
\operatorname{Viab}_F(K) = \bigcap_{T > 0} V_T(K).
\]
\end{theorem}
\begin{proof}
($\subseteq$) If $x \in \operatorname{Viab}_F(K)$, any
$\gamma \in \Gamma_K(x)$ restricts to $\Gamma_K^T(x)$ for every $T$, so
$x \in V_T(K)$ for every $T$.

($\supseteq$) Suppose $x \in \bigcap_T V_T(K)$; for each $n = 1,2,3,\dots$
choose $\gamma_n \in \Gamma_K^n(x)$. As in the sufficiency proof above,
the linear growth hypothesis (H3) gives a discrete Gr\"onwall bound
making $\{\gamma_n\}_{n \geq M}$ uniformly bounded and equi-Lipschitz on
$[0,M]$ for each fixed $M$. Arzel\`a--Ascoli and a diagonal argument over
$M$ produce a subsequence converging uniformly on every compact
$[0,M]$ to a limit curve $\gamma : [0,\infty) \to X$. Closedness of $K$
gives $\gamma(t) \in K$ for all $t$; the same weak-convergence argument
via Mazur's lemma gives $\dot\gamma(t) \in F(\gamma(t))$ for a.e.\ $t$.
Hence $\gamma \in \Gamma_K(x)$ and $x \in \operatorname{Viab}_F(K)$.
\end{proof}

\begin{remark}
This is the genuine sense in which finite-horizon objects determine the
infinite-horizon viability kernel, and it is the fact that licenses
building continuation volume and continuation distance out of
finite-horizon reachable sets, as the appendix already commits to doing:
the finite-horizon approximation is not a separate, weaker theory but a
nested sequence whose intersection recovers the object of real interest,
under exactly the hypotheses already in force throughout this book.
\end{remark}

\section*{Capture basins}

\begin{definition}[Capture basin]
For a target $C \subseteq K$ and horizon $T > 0$, define
\[
\Gamma^T_{K,C}(x) = \left\{ \gamma : [0,T] \to K \;\middle|\;
\begin{array}{l}
\gamma(0) = x,\ \dot\gamma(t) \in F(\gamma(t)) \text{ a.e.}, \\
\exists\, t^\ast \in [0,T] \text{ with } \gamma(t^\ast) \in C \text{ and } \gamma(t) \in K\ \forall t \leq t^\ast
\end{array}
\right\},
\]
and $\operatorname{Capt}_F^T(K,C) = \{x \in K : \Gamma^T_{K,C}(x) \neq \varnothing\}$.
\end{definition}

\begin{remark}
This is Aubin's own formulation, allowing capture at any time
$t^\ast \leq T$ rather than requiring it exactly at the endpoint; it is
slightly more general than the simplified fixed-endpoint version used
provisionally in the appendix, which is recovered as the special case
in which capture is required to occur exactly at $t^\ast = T$. This
is rule (a), the prescribed-final-time rule, among the three
rules distinguished in Aubin's treatment of dynamic cooperative games;
rules (b) and (c) there correspond to different choices of when capture
is permitted to count.
\end{remark}

\begin{proposition}[Monotonicity in the horizon]
$\operatorname{Capt}_F^{T_1}(K,C) \subseteq \operatorname{Capt}_F^{T_2}(K,C)$
whenever $T_1 \leq T_2$.
\end{proposition}
\begin{proof}
If $\gamma \in \Gamma^{T_1}_{K,C}(x)$ witnesses capture at some
$t^\ast \leq T_1$, the same $\gamma$, viewed on the longer domain
$[0,T_2]$ by any admissible extension past $T_1$ (which exists, since
$\gamma(t^\ast) \in C \subseteq K$ and the trajectory need not continue
past $t^\ast$ to satisfy the defining condition), already witnesses
$t^\ast \leq T_1 \leq T_2$, so $x \in \operatorname{Capt}_F^{T_2}(K,C)$.
\end{proof}

\begin{remark}[The degenerate case $C = K$]
When $C = K$, every $x \in K$ trivially satisfies the capture condition
at $t^\ast = 0$, since $\gamma(0) = x \in K = C$ by hypothesis. Hence
$\operatorname{Capt}_F^T(K,K) = K$ for every $T > 0$, regardless of
whether $K$ itself is viable. This is not the interesting connection
between capture basins and the viability kernel; that connection is
supplied by the nested-intersection theorem above, using proper
finite-horizon persistence sets $V_T(K)$ rather than a capture basin
with a degenerate target. Flagging this explicitly is worth the space,
since the superficially similar-sounding claim, that capture basins
with $C = K$ reduce to the viability kernel, is false as stated and
should not be repeated uncritically.
\end{remark}

\section*{Regulation maps}

\begin{definition}[Regulation map]
For $x \in \operatorname{Viab}_F(K)$, the regulation map is
\[
\operatorname{Reg}_K(x) = F(x) \cap T_K(x),
\]
the set of admissible velocities at $x$ that keep the system tangent to
$K$.
\end{definition}

\begin{remark}
$\operatorname{Reg}_K(x)$ is exactly the continuation operator of the
opening chapter, differentiated: it names, at each viable state, which
velocities are available for the next instant, rather than which whole
trajectories are available from here on. A continuous selection
$x \mapsto v(x) \in \operatorname{Reg}_K(x)$, where one exists, generates
an actual trajectory by solving $\dot\gamma(t) = v(\gamma(t))$; this is
the mechanism by which the existence statement of the Viability Theorem
is turned into an actual rule for staying viable, rather than a bare
guarantee that some rule exists. The next chapter takes up the case
where $\operatorname{Reg}_K(x)$ contains more than one admissible
velocity and a further criterion is needed to choose among them. This is the
point at which Aubin's own dynamic core, and this book's account of
selection, begin.
\end{remark}

\chapter{The Selection Already Inside Viability}

A ship's captain who has confirmed that some route home exists has
confirmed very little. There may be a dozen routes, and most captains
worth the title do not stop at the first one that avoids the rocks;
they ask which route survives the worst weather the season could
plausibly throw at them, and they plot that one. The difference between
those two captains is not a difference in what they know about the
sea. It is a difference in what they are willing to leave to chance
once they know a safe passage is out there somewhere.

This is worth dwelling on because it is easy to hear ``does a way
forward exist'' and ``which way forward should I take'' as two stages
of the same question, the second simply refining the first. They are
not quite that. Existence is a yes-or-no fact about a whole space of
possibilities; selection requires a standard for comparing the members
of that space against one another, and no amount of sharpening the
existence question will produce that standard on its own. What makes
the sea captain's practice worth studying formally is that the
standard she actually uses (survive the worst plausible weather) is
not an ad hoc addition bolted onto navigation from outside. It falls
directly out of taking the existence question seriously under
uncertainty, once uncertainty is treated as something an adversary
controls rather than something that merely happens. Follow the
existence question that far, and a selection criterion arrives
whether it was asked for or not.

\section*{From existence to selection}

Recall the regulation map of the previous chapter,
$\operatorname{Reg}_K(x) = F(x) \cap T_K(x)$: the set of admissible
velocities at $x$ that keep a trajectory tangent to $K$. The Viability
Theorem guarantees this set is nonempty at every viable state, but says
nothing about which of its elements to use when it contains more than
one. Viability theory, as developed so far, answers only whether
$\operatorname{Reg}_K(x)$ is empty; any measurable selection from it
generates an equally valid viable trajectory, and the theory is
indifferent among them.

To ask which selection is best requires enough additional structure to
make ``best'' mean something, and the minimal such structure separates
the velocity at each state into a part the system controls and a part
it does not.

\begin{definition}[Controlled dynamical game]
Let $P : X \rightrightarrows U$ and $Q : X \rightrightarrows V$ be
set-valued maps (the available controls, or allotments, and the
available perturbations, or tyches, at each state), and let
$f : X \times U \times V \to X$. A feedback is a map
$u_e : X \to U$ with $u_e(x) \in P(x)$ for all $x$. For a feedback
$u_e$, the trajectory set is
\[
C_{u_e}(x) = \left\{ (\gamma, v(\cdot)) \;\middle|\;
\dot\gamma(t) = f(\gamma(t), u_e(\gamma(t)), v(t)),\ v(t) \in Q(\gamma(t)),\ \gamma(0) = x \right\}.
\]
\end{definition}

\begin{proposition}[Uncontrolled viability is the trivial case]
If $P(x) = \{\star\}$ is a singleton for every $x$ (no genuine control
freedom), then setting $F(x) := \{f(x,\star,v) : v \in Q(x)\}$, the
trajectory set $C_{u_e}(x)$ under the unique feedback $u_e \equiv \star$
coincides with $\Gamma_{K}(x)$ of Chapter~1 for $K = X$: a curve
$\gamma$ appears in $C_{u_e}(x)$ (via some choice of $v(\cdot)$) if and
only if $\gamma \in \Gamma_X(x)$ under $F$.
\end{proposition}
\begin{proof}
If $(\gamma, v(\cdot)) \in C_{u_e}(x)$, then
$\dot\gamma(t) = f(\gamma(t), \star, v(t)) \in F(\gamma(t))$ for a.e.\
$t$ by definition of $F$, and $\gamma(0) = x$, so $\gamma \in \Gamma_X(x)$.
Conversely, if $\gamma \in \Gamma_X(x)$, then for a.e.\ $t$,
$\dot\gamma(t) \in F(\gamma(t)) = \{f(\gamma(t),\star,v) : v \in Q(\gamma(t))\}$,
so a measurable selection $v(t)$ with
$\dot\gamma(t) = f(\gamma(t),\star,v(t))$ exists by standard measurable
selection theorems, giving $(\gamma, v(\cdot)) \in C_{u_e}(x)$.
\end{proof}

\begin{remark}
This is the sense in which everything proved in the previous two
chapters is a special case of what follows, not a separate theory:
viability theory is controlled dynamics with the control stripped out.
Reintroducing it is not a generalization for its own sake; it is the
minimal move needed to ask which viable trajectory a system should
follow, rather than merely whether one exists.
\end{remark}

\section*{Guaranteed valuation}

Fix a horizon $T$ and a terminal payoff $u : X \to \mathbb{R}$, to be
minimized in the worst case over perturbations. Following Aubin's
treatment of dynamic cooperative games \cite{aubin2001}, and taking the
discount term $m \equiv 0$ for simplicity, define
\[
V^\sharp(T,x) = \inf_{u_e(\cdot) \in P(\cdot)} \sup_{(\gamma,v(\cdot)) \in C_{u_e}(x)} u(\gamma(T)).
\]
$V^\sharp(T,x)$ is the smallest terminal cost the system can guarantee
against every admissible perturbation, once the best available feedback
has been chosen.

\begin{definition}[The menu and the guaranteed-optimal selections]
$\Gamma(x) := P(x)$, the full menu of controls available at $x$.
$\Gamma^\star(T,x) := \{p \in P(x) : \text{using } p \text{ at } x
\text{ achieves the value } V^\sharp(T,x) \text{ against every
admissible continuation of the worst-case perturbation}\}$, made
precise via Aubin's dynamic core: when $V^\sharp$ is differentiable,
\[
\Gamma^\star(T,x) = \left\{ p \in P(x) : \sup_{v \in Q(x)} \left[ \left\langle \frac{\partial V^\sharp}{\partial x}(T,x) - p,\, f(x,v) \right\rangle \right] \leq -\frac{\partial V^\sharp}{\partial T}(T,x) \right\}.
\]
\end{definition}

\begin{remark}[A notational warning]
$\Gamma$ is being reused here for something different in kind from
$\Gamma_K(x)$ of Chapter 1: there, $\Gamma_K(x)$ was a set of whole
trajectories; here, $\Gamma(x) := P(x)$ is a set of instantaneous
control values available at $x$. The two are related, since choosing
$p \in P(x)$ at every state along a trajectory is what generates a
member of $\Gamma_K(x)$ in the first place, but they are not the same
kind of object, and conflating them would be a mistake. The symbol is
reused rather than replaced because the pattern being tracked, a menu
$\Gamma$ narrowed by some criterion to a selected subset $\Gamma^\star$,
is the same pattern each time it appears, in this chapter, in
Chapter 8, and in Chapter 9, and giving each instance its own letter
would obscure that repetition rather than clarify it. Context (whether
the surrounding discussion is about trajectories or about
instantaneous controls) disambiguates which $\Gamma$ is meant at each
occurrence from here on.
\end{remark}

\begin{remark}
This is Aubin's own dynamic-core condition, restated here with $m\equiv
0$; that the achieving-the-value characterization and the differential
inequality characterization agree is Aubin's own result and is not
re-derived here \cite{aubin2001}. What matters for this chapter is only
that $\Gamma^\star(T,x)$ is defined as a subset of $\Gamma(x) = P(x)$ by
an extra condition beyond mere availability, so trivially
$\Gamma^\star(T,x) \subseteq \Gamma(x)$, and the question is whether
that inclusion is ever strict in a way that matters.
\end{remark}

\begin{proposition}
$\Gamma^\star(T,x) \subseteq \Gamma(x)$ for every $T, x$.
\end{proposition}
\begin{proof}
Immediate: $\Gamma^\star(T,x)$ is defined as $\{p \in P(x) : \dots\}$, a
subset of $P(x) = \Gamma(x)$ by construction.
\end{proof}

\section*{A worked example}

The inclusion above is trivial as a set-theoretic fact; what is not
trivial is that it is often strict, and that the gap it leaves out
represents controls that are viable in the sense of the previous
chapter yet indefensible in the sense of this one.

\begin{example}
Let $X = \mathbb{R}$, $P(x) = \{-1, +1\}$ for all $x$ (two available
constant controls), $Q(x) = [-\tfrac12, \tfrac12]$ (a bounded
perturbation), and $f(x,u,v) = u + v$. Take terminal payoff
$u(x) = x^2$ and horizon $T > 0$, with the system started at $x = 0$.

Consider two feedbacks: $u_e^{(1)}(x) \equiv +1$ for all $x$, and
$u_e^{(2)}(x) = -\operatorname{sign}(x)$ for $x \neq 0$ (with the
value at $x=0$ resolved by the Filippov sliding-mode convention
\cite{filippov1988,utkin1992}, as is standard for discontinuous
feedbacks of this type).
\end{example}

\begin{proposition}
Under $u_e^{(1)}$, $\sup_{(\gamma,v(\cdot))} \gamma(T)^2 = \tfrac{9}{4}T^2$.
Under $u_e^{(2)}$, $\sup_{(\gamma,v(\cdot))} \gamma(T)^2 = 0$. Hence
$V^\sharp(T,0) = 0$, $u_e^{(2)} \in \Gamma^\star(T,0)$, and
$u_e^{(1)} \notin \Gamma^\star(T,0)$, even though both are available
selections: $u_e^{(1)}(0), u_e^{(2)}(0) \in P(0) = \Gamma(0)$.
\end{proposition}
\begin{proof}
Under $u_e^{(1)}$, $\dot\gamma(t) = 1 + v(t)$ with $v(t) \in [-\tfrac12,
\tfrac12]$ chosen by the adversary independently at each instant. Since
$\gamma(t) = \int_0^t (1 + v(s))\, ds$ and the integrand is
pointwise maximized by $v(s) = \tfrac12$, the supremum of $\gamma(T)$
over admissible $v(\cdot)$ is attained at $v \equiv \tfrac12$, giving
$\gamma(T) = \tfrac32 T > 0$, hence $\sup \gamma(T)^2 = \tfrac94 T^2$.

Under $u_e^{(2)}$, consider $W(t) = \gamma(t)^2$. Wherever
$\gamma(t) \neq 0$,
\[
\dot W(t) = 2\gamma(t)\dot\gamma(t) = 2\gamma(t)\big(-\operatorname{sign}(\gamma(t)) + v(t)\big)
= -2|\gamma(t)| + 2\gamma(t)v(t) \leq -2|\gamma(t)| + |\gamma(t)| = -|\gamma(t)| \leq 0,
\]
using $|v(t)| \leq \tfrac12$. So $W$ is nonincreasing whenever
$\gamma(t) \neq 0$. Starting from $\gamma(0) = 0$, i.e.\ $W(0) = 0$,
and $W \geq 0$ always, any departure from $W = 0$ would require
$\dot W > 0$ at some point with $\gamma \neq 0$, which the inequality
above forbids; the Filippov sliding solution at $\gamma = 0$
correspondingly keeps $\gamma(t) \equiv 0$ for all $t$, for every
admissible $v(\cdot)$ \cite{filippov1988}. Hence $\gamma(T) \equiv 0$
regardless of the adversary's choice, so $\sup \gamma(T)^2 = 0$.

Since $\gamma(T)^2 \geq 0$ always, $0$ is the global lower bound for
the guaranteed value, and $u_e^{(2)}$ achieves it, so
$V^\sharp(T,0) = 0$ and $u_e^{(2)} \in \Gamma^\star(T,0)$. Since
$u_e^{(1)}$ achieves only $\tfrac94T^2 > 0 = V^\sharp(T,0)$, it fails to
achieve the guaranteed value and so $u_e^{(1)} \notin \Gamma^\star(T,0)$.
\end{proof}

\begin{remark}
Both feedbacks are viable in the sense of Chapter 3: both keep the
trajectory within $X = K$ trivially, since there is no constraint set
here beyond the whole line. Viability theory alone cannot distinguish
them: $\operatorname{Reg}_K(x) = \{-1,+1\}$ regardless of $x$, and
both elements are equally admissible from that vantage point. It is
only once a worst-case standard is imposed that one of the two
becomes visibly indefensible: $u_e^{(1)}$ lets the adversary drive the
terminal state arbitrarily far from the target as $T$ grows, while
$u_e^{(2)}$ holds it exactly at the target regardless of what the
adversary does. This is the concrete content behind the abstract claim
that $\Gamma^\star(T,x) \subsetneq \Gamma(x)$ can be a strict and
consequential inclusion, not merely a strict one in principle.
\end{remark}

\section*{What this does and does not settle}

The example makes the selection performed by $V^\sharp$ concrete, but
it should not be read as settling the question of \emph{which}
selection criterion a continuation should be judged by in general.
$V^\sharp$ selects for the best guaranteed outcome against an adversary
that is permitted to see the system's feedback and respond
optimally against it: a strong and specific standard, appropriate
when the source of uncertainty really is adversarial or must be
planned against as if it were. Later chapters introduce criteria
answering to recoverability, to the preservation of distinctions, and
to the retention of future options. None of them reduce to worst-case
guarantee against an adversary, and none of them is more fundamental
than it. What this chapter has shown is only that the general pattern,
a menu $\Gamma$ narrowed by some criterion to $\Gamma^\star$, is
not a novelty this book is introducing from outside viability theory.
It is already how Aubin's own construction behaves, under a criterion
this book did not have to invent. The chapters that follow reuse the
pattern and change only the criterion.

\chapter{Continuation Operators}

Two patches of skin can look, under a microscope, identical, and still
disagree completely about what they are capable of becoming. One is
unmarked tissue with the ordinary regenerative range of the organism
around it; the other is a scar, chemically and structurally
indistinguishable from ordinary skin in a great many respects, and yet
closed off from futures the first patch still has open to it. Nothing
about the scar's present molecular configuration announces this
difference to a description that only looks at its current state. The
difference is legible only in how each patch got to where it is.

The same pattern turns up anywhere a snapshot undersells its subject. A
codebase with a thousand small, well-tested commits behind it and a
codebase that reached the identical set of files through one enormous,
undocumented rewrite are not equally easy to modify tomorrow, even if
`git diff` between them today reports nothing. A position on a chess
board carries different practical continuations depending on how much
time each player has left on their clock, information the board itself
does not encode. In each case, the object this book has been calling
the state is not quite doing its job: two systems can occupy the same
state and still face different futures, because the future available
to a system is sometimes a fact about its path, not only about its
present position.

This is a genuine complication for everything built so far, since every
construction up to this point (the viability kernel, the guaranteed
value, the regulation map) has been written as a function of the
current state alone. The complication has a name in the sciences that
study it directly: memory, in the broad sense that includes a scar,
a commit history, and a chess clock alike. What follows is not a
solution to memory in general, which would be too much to ask of one
chapter, but a precise account of when memory can be tamed, folded
into a slightly larger notion of state without losing anything, and
when it genuinely cannot, so that the rest of the book knows which kind
of problem it is dealing with at any given point.

\section*{The trivial embedding}

\begin{definition}[History]
For a trajectory $\gamma : [0,\infty) \to X$, the history of $\gamma$
up to time $t$ is $H_t := \gamma|_{[0,t]}$, an element of
$\mathcal{H}_t := C([0,t], X)$, the space of continuous $X$-valued
paths on $[0,t]$.
\end{definition}

\begin{definition}[History-dependent continuation operator]
A continuation problem is history-dependent if the admissible
continuations from $x_t := \gamma(t)$ depend on $H_t$ and not only on
$x_t$: there is a map $\Gamma(\,\cdot \mid \cdot\,)$ with
$\Gamma(x_t \mid H_t) \subseteq \Gamma(x_t \mid H_t')$ failing in
general for $H_t \neq H_t'$ even when both end at $\gamma(t) = x_t$.
\end{definition}

\begin{proposition}[Universal embedding]
Every history-dependent continuation problem embeds into a state-only
(Markovian) continuation problem on the augmented space
$\tilde X := X \times \bigcup_{t \geq 0} \mathcal{H}_t$, via
$\tilde x_t := (x_t, H_t)$, whose evolution is deterministic given
$\gamma$: $\tilde x_{t+dt} = (x_{t+dt}, H_t \frown [\gamma]_{[t,t+dt]})$,
extending $H_t$ by the newly traversed segment.
\end{proposition}
\begin{proof}
By construction, $H_t$ is exactly $\gamma|_{[0,t]}$, so $H_{t+dt}$ is
determined without ambiguity by $H_t$ together with the segment of
$\gamma$ traversed on $[t, t+dt]$; no information beyond
$(x_t, H_t)$ and the dynamics generating the next increment is needed
to determine $\tilde x_{t+dt}$. Hence $\Gamma(x_t \mid H_t)$, as a
function of $H_t$ alone (with $x_t$ recoverable as the endpoint of
$H_t$), is a well-defined function of $\tilde x_t$ alone, which is the
definition of a state-only continuation operator on $\tilde X$.
\end{proof}

\begin{remark}
This proposition resolves half of the open question left at the end of
the appendix, but only in the weakest possible sense, and it is worth
being honest about how weak. $\tilde X$ is generally
infinite-dimensional ($\mathcal{H}_t$ is a function space), and the
embedding does no explanatory work: it says that history-dependence can
always be redescribed as state-dependence provided the state is allowed
to be the entire history, which is close to restating the problem
rather than solving it. The substantive question is not whether some
augmentation exists, but whether a \emph{small} one does: whether the
dependence on $H_t$ factors through some much smaller summary of it.
\end{remark}

\section*{Sufficient statistics and genuine reduction}

\begin{definition}[Sufficient statistic for continuation]
A map $\varphi : \bigcup_t \mathcal{H}_t \to M$, for some space $M$, is
a sufficient statistic for a history-dependent continuation problem if
$\varphi(H_t) = \varphi(H_t')$ implies
$\Gamma(x_t \mid H_t) = \Gamma(x_t \mid H_t')$ whenever both histories
end at the same $x_t$. $\varphi$ is autonomous if there is a map
$g : X \times M \to M$ (or, in continuous time, a dynamics on $M$
driven by $x_t$) such that $\varphi(H_{t+dt})$ is determined by
$x_t$, $\varphi(H_t)$, and the dynamics alone, without reference to any
other part of $H_t$.
\end{definition}

\begin{proposition}[Autonomous sufficient statistics give genuine reduction]
If $\varphi$ is an autonomous sufficient statistic taking values in a
finite-dimensional $M$, then the augmented state $(x_t, \varphi(H_t)) \in
X \times M$ carries a genuine Markovian continuation operator: there is
a well-defined $\hat\Gamma$ on $X \times M$ with
$\hat\Gamma(x_t, \varphi(H_t)) = \Gamma(x_t \mid H_t)$ for every history
$H_t$, and the evolution of $(x_t, \varphi(H_t))$ is itself Markovian
on $X \times M$.
\end{proposition}
\begin{proof}
Well-definedness of $\hat\Gamma(x,m) := \Gamma(x \mid H)$ for any
history $H$ with endpoint $x$ and $\varphi(H) = m$ is exactly the
sufficiency condition: any two such histories give the same value of
$\Gamma$, so the choice of representative $H$ does not matter. Markovian
evolution of $(x_t, \varphi(H_t))$ is exactly the autonomy condition:
$\varphi(H_{t+dt})$ is determined by $(x_t, \varphi(H_t))$ and the
dynamics alone, so the pair evolves without reference to any
information beyond itself. Since $\dim(X \times M) < \infty$ by
hypothesis, this is a genuine finite-dimensional reduction, not the
trivial embedding of the previous section.
\end{proof}

\begin{example}[A natural sufficient statistic]
Let $X = \mathbb{R}$, and suppose $\Gamma(x_t \mid H_t)$ depends on
$H_t$ only through the running maximum $\varphi(H_t) := \max_{s \leq t} \gamma(s)$,
as occurs, for instance, whenever the continuations available to a
system depend on the highest level it has ever reached, and not on the
particular path taken to get there or back. Then $\varphi$ is a
sufficient statistic by hypothesis, and it is autonomous: given
$\dot\gamma(t) = v(t)$,
\[
\frac{d}{dt}\varphi(H_t) = \begin{cases} v(t) & \text{if } \gamma(t) = \varphi(H_t) \text{ and } v(t) \geq 0, \\ 0 & \text{otherwise,} \end{cases}
\]
a rule depending only on $x_t = \gamma(t)$, $\varphi(H_t)$, and the
current velocity, not on any earlier part of $H_t$. By the previous
proposition, $(x_t, \varphi(H_t)) \in \mathbb{R}^2$ carries the full
continuation operator, a reduction from an infinite-dimensional history
to two real numbers.
\end{example}

\begin{remark}
This is a standard construction in stochastic control and mathematical
finance, where lookback-style payoffs depending on a running maximum or
minimum are handled by exactly this augmentation \cite{shreve2004}; it
is recorded here because it is the cleanest illustration available of
what a genuine reduction looks like, as opposed to the trivial embedding
above.
\end{remark}

\section*{A limit to reduction}

The examples above might suggest that a finite sufficient statistic can
always be found if one is only clever enough to look for it. This is
false, and it is worth having a clean proof that it is false, so that
the search for a sufficient statistic in any particular application is
understood as a real question with a real chance of failing, not a
formality.

\begin{proposition}[Existence of continuation problems with no finite-range sufficient statistic]
There exist history-dependent continuation problems admitting no
sufficient statistic $\varphi$ taking only finitely many values:
that is, no sufficient statistic of finite \emph{range}, in the sense
of finite cardinality of $\varphi$'s image.
\end{proposition}
\begin{proof}
Fix discrete time $t = 0, 1, 2, \dots$, $X = \{0,1\}$, and define
$\Gamma(x_t \mid H_t) := \{H_t\}$: the continuation operator simply
returns the history itself as its own (singleton) value. Suppose
$\varphi$ is sufficient with range $M$ of fixed finite cardinality.
Sufficiency requires that $\varphi(H_t) = \varphi(H_t')$ implies
$\Gamma(x_t \mid H_t) = \Gamma(x_t \mid H_t')$, i.e.\ $\{H_t\} = \{H_t'\}$,
i.e.\ $H_t = H_t'$. So $\varphi$ must be injective on
$\mathcal{H}_t = \{0,1\}^t$. But $|\mathcal{H}_t| = 2^t$, which exceeds
$|M|$ once $t > \log_2 |M|$, so no map from a $2^t$-element set to $M$
can be injective, contradicting sufficiency. Hence no $\varphi$ with a
fixed finite range suffices for every $t$.
\end{proof}

\begin{remark}[This does not rule out finite-dimensional statistics]
The proof above uses only cardinality: a set of fixed finite size
cannot injectively absorb an unboundedly growing set. It says nothing
about statistics valued in a finite-\emph{dimensional} but
infinite-\emph{cardinality} space such as $M = \mathbb{R}^k$, which
remains entirely uninjured by this argument. $|\mathbb{R}^k| =
|\mathbb{R}|$ regardless of $k$, comfortably large enough to injectively
absorb $\mathcal{H}_t$ for every $t$, so the proposition above cannot be
strengthened to rule out real-valued statistics by this method. Whether
some dimensional obstruction exists (some topological or
measure-theoretic reason a finite-dimensional statistic must fail even
when a merely finite-range one is already ruled out) is a harder
question, requiring different tools than a counting argument, and is
left open here rather than claimed.
\end{remark}

\begin{remark}
This construction is adversarial by design: $\Gamma$ was built to
depend on literally the entire history, with nothing discarded, so of
course nothing smaller than the entire history can summarize it. Most
continuation problems that arise from actual dynamics, rather than
being handed to us as an arbitrary function of history, are not like
this; the running-maximum example above is closer to the typical case.
But the existence of even one such problem is enough to establish that
sufficiency is a substantive constraint, not something guaranteed by
the mere fact that $X$ and the class of admissible $\Gamma$ are both
``nice.'' Whether a given system, Levin's bioelectric memory among
them, admits a finite or low-dimensional sufficient statistic is
therefore an empirical and structural question about that system, to
be settled case by case, not a question this chapter's machinery
answers in general.
\end{remark}

\section*{Where this leaves the state}

For every construction later in this book, one of two things is true.
Either a sufficient statistic has been found or is assumed to exist, in
which case $x$ silently means the augmented pair $(x, \varphi(H_t))$
from here on, and every later definition, volume, distance,
admissibility, inherits that augmentation without needing to be
rewritten; or no such statistic is assumed, in which case $\Gamma$
should be read as history-dependent throughout, and any later
proposition proved only for the state-only case should be understood as
provisional pending an argument for why the relevant sufficient
statistic exists in the application at hand. This book does not commit
to one option over the other in general. What it commits to is the
discipline of noticing, in each later chapter, which option is in
force.

\chapter{Continuation Spaces}

Two trailheads can offer a hiker the same number of marked paths and
still not offer the same amount of freedom. At the first, five trails
fan out from the junction at wide, distinct angles, and a hiker who
changes her mind after the first hundred yards is still choosing among
five genuinely different directions. At the second, five trails leave
the junction too, but four of them run within sight of each other for
the first mile before finally diverging; a hiker who changes her mind
early has, in every practical sense, only two real options, not five.
Counting the trails told us nothing about this difference. What told
us something was a measure of how much room the trails actually
occupy, and a measure of how close together they run.

This chapter is about building those two measures, room and
closeness, for the continuation spaces this book has been
constructing since the opening chapter, and about being honest, at
each step, about which of the two can currently be built on solid
ground and which cannot yet. Both turn out to depend on a decision
made early, in the very first chapter: whether to measure the whole
infinite future at once, or to measure it one finite horizon at a time
and treat the infinite case as a limit. The choice matters more here
than it has anywhere so far, because it is the difference between a
measure that can be computed and one that, it turns out, cannot be
built at all in the most natural-seeming way.

\section*{Why the naive volume fails}

The most direct way to measure how much room a continuation space
occupies is to put a measure directly on the trajectory space itself.

\begin{definition}[Naive continuation volume]
$V_\Gamma(x) := \mu(\Gamma_K(x))$, for some measure $\mu$ on the
infinite-horizon trajectory space $\Gamma_K(x) \subset C([0,\infty), K)$.
\end{definition}

The natural candidate for $\mu$, wherever the dynamics involves genuine
uncertainty, is a Wiener-type measure: the distribution induced on path
space by a stochastic differential equation whose drift is compatible
with $F$. The following proposition shows this candidate fails outright,
not approximately, for the trajectory spaces this book has been using.

\begin{proposition}[Wiener-type measures vanish on continuation spaces]
Let $\mu$ be Wiener measure on $C([0,T],X)$, or any measure absolutely
continuous with respect to it. Then $\mu(\Gamma_K^T(x)) = 0$.
\end{proposition}
\begin{proof}[Proof sketch]
By construction, every $\gamma \in \Gamma_K^T(x)$ satisfies
$\dot\gamma(t) \in F(\gamma(t))$ for a.e.\ $t$, which requires $\gamma$
to be absolutely continuous; absolutely continuous functions are of
bounded variation on $[0,T]$. But a $\mu$-typical path (a.s.\ with
respect to Wiener measure) is nowhere differentiable
\cite{paleywienerzygmund1933} and has unbounded variation on every
subinterval, hence is almost surely \emph{not} absolutely continuous.
So the set of absolutely continuous paths has $\mu$-measure zero, and
$\Gamma_K^T(x)$, being a subset of that set, has $\mu$-measure zero as
well. This argument is elementary once the nowhere-differentiability
and unbounded-variation properties of Brownian-type paths are granted,
but those properties are themselves a nontrivial classical result
\cite{paleywienerzygmund1933} not reproved here, which is why this is
labeled a sketch rather than a self-contained proof, in keeping with
how Chapter 3 labels its own reliance on comparably deep cited results.
\end{proof}

\begin{remark}
This makes precise, and turns into a short clean proof, what the
appendix flagged only as a known difficulty: it is not that Wiener-type
measures are an awkward or ungraded way to measure continuation spaces;
it is that they assign the continuation space zero measure outright,
for the elementary reason that Brownian-type randomness and the
bounded-variation trajectories this book studies are, in a precise
sense, mutually singular. No amount of refining the choice of Wiener
measure repairs this, since the argument uses nothing about the
particular drift or diffusion coefficients, only that $\mu$-typical
paths are nowhere differentiable. $V_\Gamma$ as originally proposed is
therefore not merely difficult to compute; it is identically zero
whenever built this way, which is worse than being merely graded
poorly.
\end{remark}

\section*{Volume on finite horizons}

The finite-horizon alternative recommended in the appendix survives
exactly because it does not attempt to measure a space of paths at all;
it measures a space of endpoints.

\begin{definition}[Reachable set and continuation volume]
$R_T(x) := \{\gamma(T) : \gamma \in \Gamma_K^T(x)\} \subseteq K$.
$\operatorname{Vol}_T(x) := \lambda(R_T(x))$, where $\lambda$ is
Lebesgue measure on $X = \mathbb{R}^n$ (or Hausdorff measure of the
appropriate dimension, if $R_T(x)$ is lower-dimensional).
\end{definition}

\begin{remark}
Notation: this is distinguished from $V_T(K)$ of the previous chapter,
which denoted a \emph{subset of states} admitting a $T$-length
trajectory. $\operatorname{Vol}_T(x)$ is a number, the size of the
reachable set from one fixed starting point.
\end{remark}

\begin{proposition}[Well-posedness]
Under the standing hypotheses of Chapter 3, $R_T(x)$ is compact for
every $x \in K$ and $T > 0$, and hence Lebesgue measurable; so
$\operatorname{Vol}_T(x)$ is well-defined.
\end{proposition}
\begin{proof}
This is the classical compactness theorem for reachable sets of
differential inclusions with compact convex values, upper
semicontinuous right-hand side, and linear growth
\cite{aubincellina1984,filippov1988}: under these hypotheses, the map
$x \mapsto R_T(x)$ is itself compact-valued and upper semicontinuous.
Compact subsets of $\mathbb{R}^n$ are closed and bounded, hence Borel,
hence Lebesgue measurable.
\end{proof}

\begin{proposition}[Monotonicity in the constraint set]
If $K_1 \subseteq K_2$, then $\operatorname{Vol}_T^{K_1}(x) \leq \operatorname{Vol}_T^{K_2}(x)$.
\end{proposition}
\begin{proof}
This is Proposition 5 of Chapter 3, applied at the fixed starting point
$x$: $R_T^{K_1}(x) \subseteq R_T^{K_2}(x)$ by the same inclusion
argument, and monotonicity of Lebesgue measure under inclusion gives
the result.
\end{proof}

\begin{remark}[Volume is not monotone in the horizon]
It might be expected, by analogy with the previous proposition, that
$\operatorname{Vol}_T(x)$ grows with $T$: more time, more room. This is
false in general, and it is worth a clean counterexample rather than
leaving the failure to intuition.
\end{remark}

\begin{example}
Let $X = \{0,1,2\}$ (discrete time, discrete state space), with
one-step reachability $\Phi(0) = \{1,2\}$, $\Phi(1) = \{0\}$,
$\Phi(2) = \{0\}$, and let $\operatorname{Vol}$ be counting measure.
Starting from $x_0 = 0$: $R_0(0) = \{0\}$, so
$\operatorname{Vol}_0(0) = 1$. $R_1(0) = \Phi(0) = \{1,2\}$, so
$\operatorname{Vol}_1(0) = 2$. $R_2(0) = \Phi(1) \cup \Phi(2) = \{0\} \cup \{0\} = \{0\}$,
so $\operatorname{Vol}_2(0) = 1 < \operatorname{Vol}_1(0)$: volume
strictly \emph{decreased} from $T=1$ to $T=2$, because the two branches
reached at $T=1$ collapse back onto a single state at $T=2$.
$R_3(0) = \Phi(0) = \{1,2\}$ again, so the sequence oscillates:
$1, 2, 1, 2, \dots$
\end{example}

\begin{remark}
This is not a pathology to be engineered away; it is a direct
discrete-time descendant of the Poincar\'e recurrence discussed in
Chapter 2. Reachable sets can fold back on themselves, and when they
do, their volume can shrink even though every individual trajectory
persists. Monotonicity under constraint tightening (Proposition 5,
reused above) is a fact about set inclusion and holds unconditionally;
monotonicity in time is a fact about dynamics and holds only when the
dynamics happens to cooperate.
\end{remark}

\section*{Continuation distance}

\begin{definition}[Continuation distance]
$d_\Gamma(x,y) := d_H(R_T(x), R_T(y))$, the Hausdorff distance between
the finite-horizon reachable sets, on the ambient metric of $X$.
\end{definition}

\begin{proposition}
$d_\Gamma$ is a pseudometric on $K$: symmetric, satisfying the triangle
inequality, and satisfying $d_\Gamma(x,x) = 0$; it is a genuine metric
on the quotient identifying states with equal reachable sets.
\end{proposition}
\begin{proof}
Symmetry and $d_\Gamma(x,x)=0$ are immediate from the definition of
Hausdorff distance and $R_T(x) = R_T(x)$. For the triangle inequality,
recall $d_H(A,B) = \max\{\sup_{a \in A} \operatorname{dist}(a,B),\
\sup_{b \in B} \operatorname{dist}(b,A)\}$; for compact $A,B,C$, every
point of $A$ is within $d_H(A,B)$ of some point of $B$, which is in
turn within $d_H(B,C)$ of some point of $C$, so every point of $A$ is
within $d_H(A,B) + d_H(B,C)$ of $C$, and symmetrically for $C$ against
$A$, giving $d_H(A,C) \leq d_H(A,B) + d_H(B,C)$. Degeneracy on the
quotient (rather than on $K$ itself) follows because $d_H(A,B) = 0$
implies $A = B$ only when both are closed, which $R_T(x)$ is by the
compactness proposition above, so $d_\Gamma(x,y) = 0$ implies
$R_T(x) = R_T(y)$ exactly, not merely approximately.
\end{proof}

\begin{remark}[Which distance is the right one]
Hausdorff distance measures worst-case set overlap: it is large exactly
when some part of one reachable set is far from every part of the
other. This is not the only reasonable notion, and the appendix's open
problem on this point is not resolved here. Three alternatives deserve
naming, since each answers a different question. If the reachable sets
carried their own intrinsic geometry that differed from state to
state, useful if $R_T(x)$'s internal structure (not just its embedding
in $X$) is what later chapters care about, Gromov--Hausdorff distance
would be the natural choice. Suppose instead $R_T(x)$ carried a
density: the distribution of endpoints under a natural measure on
admissible controls, say, rather than being treated as a bare set.
Then Wasserstein distance becomes appropriate, exactly when volume
itself is upgraded from Lebesgue measure of a set to a genuine
probability-weighted quantity. And if the point of the distance is to
measure how hard it is to convert one continuation space into another,
rather than how far apart they are as sets, what is wanted is a
reachability-based metric: the minimal control effort needed to move
the system from a state whose reachable set resembles $R_T(x)$ to one
resembling $R_T(y)$. Which of these is wanted depends on which later
construction is doing the asking, and this book does not yet have
enough worked examples to settle the question in general.
\end{remark}

\section*{Continuation density}

\begin{definition}[Continuation density]
Where $\operatorname{Vol}_T$ is differentiable at $x$ in direction $u$,
\[
\rho_\Gamma(x;u) := D\operatorname{Vol}_T(x)[u] = \lim_{h \to 0} \frac{\operatorname{Vol}_T(x+hu) - \operatorname{Vol}_T(x)}{h}.
\]
\end{definition}

\begin{example}
Let $X = \mathbb{R}^n$, $K = X$, and $F(x) = \overline{B}(0, r(x))$, the
closed ball of radius $r(x)$ around the origin, for some smooth
$r : X \to (0,\infty)$. Then $R_T(x)$ is (to leading order in $T$, for
small $T$) approximately the ball $\overline{B}(x, r(x)T)$, giving
\[
\operatorname{Vol}_T(x) \approx c_n\, r(x)^n T^n
\]
for the volume constant $c_n$ of the unit ball in $\mathbb{R}^n$, and
hence
\[
\rho_\Gamma(x;u) \approx c_n n\, r(x)^{n-1} T^n\, \nabla r(x) \cdot u,
\]
a directional derivative computable directly from the smooth dependence
of the velocity bound $r$ on position.
\end{example}

\begin{remark}
This example is deliberately the simplest nondegenerate case,
constant-direction growth from a velocity ball of state-dependent
radius, and is meant only to show that $\rho_\Gamma$ is a genuinely
computable quantity in at least one class of examples, not to claim
that $\operatorname{Vol}_T$ is differentiable in general. Differentiability
of $\operatorname{Vol}_T$ depends on regularity of the reachable-set
boundary that has not been established for arbitrary $F$ and is not
assumed here.
\end{remark}

\section*{Curvature}

No definition is proposed. Curvature ordinarily presupposes a
Riemannian or at least a well-behaved metric structure: geodesics,
a notion of parallel transport, a second-order comparison between
nearby distances, and none of that has been earned yet. Even
continuation distance, established above only as a Hausdorff
pseudometric on finite-horizon reachable sets, has not been shown to
admit anything like a smooth geodesic structure on the space of such
sets. Introducing curvature now would mean inventing a definition to
fit a word rather than discovering what the word could mean once volume
and distance have been tested against enough examples to show what
shape they actually have. That test is future work, not owed by this
chapter.

\section*{What this geometry does not yet say}

Volume and distance, as built here, are answers to purely quantitative
questions: how much room, how far apart. Neither says anything about
which parts of a continuation space matter more than others: whether
a given cubic unit of reachable space represents options that are
redundant with each other or options that are genuinely distinct, in
whatever sense the system's own distinctions make precise. That
question is deferred deliberately to the next chapter, where volume and
distance stop being the whole story and start being raw material for
something that cares about content, not merely extent.

\chapter{Distinguishability and Repair}

A damaged manuscript can be restored in two quite different senses.
A restorer can fill the torn-out passage with ink of the right color,
in a hand that matches the surrounding script, producing a page that
is once again whole, legible, and pleasant to look at. Or a restorer
can recover, from some other witness (a copyist's letter, a
quotation in a later work, a watermark that dates the paper) what
the passage actually said, and write that back in. Both restorers hand
back an intact page. Only one of them has given back the thing that
was lost. The first kind of repair is real, and sometimes it is all
that is available, but it should never be confused with the second
kind, and a theory of repair that cannot tell them apart has missed
the part of the problem that made repair worth caring about in the
first place.

Everything this book has built so far can already say whether a
manuscript (or a body, or a system of any kind) can be returned
to some designated condition. That is a question about existence, and
it has been fully answered since Chapter 3. What it cannot yet say is
whether the return preserves the thing that mattered, as opposed to
merely producing something that satisfies the letter of the target
condition while quietly discarding what made the original original.
Saying that requires a formal account of what it means for two
outcomes to differ in the first place, an account of distinction, and
building one, carefully enough that it connects to the geometry
already in hand rather than floating free of it, is the business of
this chapter.

\section*{Repair as capture}

\begin{definition}[Repair problem]
A repair target $R \subseteq X$ names a set of states counted as
recovered. The repair space of $x$ is
\[
\operatorname{Repair}_T(x) := \{ \gamma : [0,T] \to X \mid \gamma(0) = x,\ \exists\, t^\ast \leq T,\ \gamma(t^\ast) \in R,\ \dot\gamma(t) \in F(\gamma(t)) \}.
\]
\end{definition}

By Chapter 3's definition of the capture basin with $C := R$, $K := X$,
$x \in \operatorname{Capt}_F^T(X,R)$ if and only if
$\operatorname{Repair}_T(x) \neq \varnothing$: this is nothing more than
the capture basin, renamed for the application. Existence of a repair
trajectory is therefore already fully answered by machinery already in
hand.

Construction is a different question, and it is the one that actually
matters to a system trying to repair itself under conditions it does
not control. Chapter 4 distinguished the menu of available controls,
$\Gamma(x)$, from the guaranteed-optimal subset $\Gamma^\star(T,x)$
selected by worst-case performance against an adversary. The repair
problem is a direct instance of that same construction.

\begin{proposition}[Guaranteed repair is guaranteed valuation]
Let $u(x) := \operatorname{dist}(x, R)$ be the terminal payoff. Then
the guaranteed-optimal selections $\Gamma^\star(T,x)$ of Chapter 4,
computed with this payoff, are exactly the feedbacks that guarantee
repair against every admissible perturbation: $u_e \in \Gamma^\star(T,x)$
if and only if every trajectory in $C_{u_e}(x)$ reaches $R$ by time $T$
whenever $V^\sharp(T,x) = 0$, and otherwise achieves the smallest
guaranteed terminal distance to $R$ available under any feedback.
\end{proposition}
\begin{proof}
This is the definition of $\Gamma^\star(T,x)$ from Chapter 4, specialized
to the payoff $u(x) = \operatorname{dist}(x,R)$: a feedback achieves the
guaranteed value $V^\sharp(T,x) = \inf_{u_e} \sup_{v(\cdot)} \operatorname{dist}(\gamma(T), R)$
exactly when it is a guaranteed-optimal selection in the sense already
established. When $V^\sharp(T,x) = 0$, achieving it means guaranteeing
$\operatorname{dist}(\gamma(T),R) = 0$ against every perturbation, i.e.\
guaranteeing $\gamma(T) \in R$ (as $R$ is closed), which is guaranteed
repair.
\end{proof}

\begin{remark}
This is what separates a repair strategy from a repair certificate.
Knowing $x \in \operatorname{Capt}_F^T(X,R)$ certifies that repair is
possible for \emph{some} perturbation sequence; it says nothing about
what to do if the perturbation turns out to be adversarial. A feedback
in $\Gamma^\star(T,x)$ is a strategy that works regardless. Everything
this book has said about the gap between $\Gamma$ and $\Gamma^\star$ in
Chapter 4 applies here without modification. Repair theory did not
need a new selection mechanism, only the existing one pointed at a
repair target instead of an arbitrary payoff.
\end{remark}

\section*{What a distinction is}

Neither of the constructions above says anything about whether a
repaired state has preserved what an observer of the system actually
cares about. Both restorers from the opening paragraph produce a
$\gamma$ with $\gamma(T) \in R$, if $R$ is taken to mean simply
``intact page.'' Telling them apart requires naming, formally, what
it would mean for the two restorations to differ.

\begin{definition}[Distinction]
A distinction on $K$ is a partition $\mathcal{D} = \{D_i\}_{i \in I}$
of $K$ into disjoint nonempty cells. Two states $x, y \in K$ are
distinguished by $\mathcal{D}$ if they lie in different cells.
\end{definition}

This is deliberately the weakest possible formal object answering to
the word: no metric structure, no ordering, nothing beyond ``treated as
the same'' or ``treated as different.'' What makes it useful here is
that the continuation distance of the previous chapter gives a way to
ask whether the dynamics respects a distinction, without needing to add
anything further to the definition.

\begin{definition}[Distinction margin]
For a distinction $\mathcal{D}$ and $x \in K$ with cell $D(x)$, the
distinction margin of $x$ at horizon $T$ is
\[
A_{\mathcal{D}}^T(x) := \inf_{y \notin D(x)} d_\Gamma^T(x,y),
\]
using the continuation distance $d_\Gamma^T(x,y) = d_H(R_T(x), R_T(y))$
of the previous chapter.
\end{definition}

\begin{proposition}[Zero margin is exactly collapse]
$A_{\mathcal{D}}^T(x) = 0$ if and only if there exists $y \notin D(x)$
with $R_T(x) = R_T(y)$, provided the infimum defining $A_{\mathcal D}^T(x)$
is attained (in particular, whenever $K \setminus D(x)$ is compact).
\end{proposition}
\begin{proof}
If the infimum is attained at some $y_0 \notin D(x)$ with
$d_\Gamma^T(x,y_0) = 0$, then by the metric property established in
Chapter 6 (equal reachable sets give zero distance, and conversely on
the relevant quotient), $R_T(x) = R_T(y_0)$. Conversely, if
$R_T(x) = R_T(y)$ for some $y \notin D(x)$, then $d_\Gamma^T(x,y) = 0$,
so the infimum is at most $0$, and since $d_\Gamma^T \geq 0$ always,
the infimum equals $0$.
\end{proof}

\begin{remark}
This is the precise sense in which a distinction can be said to
collapse under continuation: not when two states merely look similar,
but when their entire future reachable sets, at some fixed horizon,
coincide exactly. An observer restricted to asking ``what can this
system still become'' literally cannot tell $x$ and $y$ apart once
$A_{\mathcal D}^T(x) = 0$, regardless of what distinction $\mathcal D$
declared between them a moment before. This is also, incidentally, a
natural formal cousin of the diagnostic-coordinate question: how
much information remains available to resolve a state from its
alternatives. The connection is only noted here, not
developed; diagnostic coordinates as originally posed concern
available information at a point in time rather than distance in a
space of futures, and the precise relationship between the two remains
outside this book's scope.
\end{remark}

\begin{definition}[Distinction preservation along a trajectory]
For $\gamma \in \Gamma_K^T(x)$, define
\[
A_{\mathcal D}(\gamma) := \inf_{t \in [0,T]} A_{\mathcal D}^T(\gamma(t)),
\]
the worst-case distinction margin maintained along $\gamma$.
\end{definition}

\begin{remark}
This is exactly the shape Chapter 8 will need: a functional
$A_{\mathcal D} : \Gamma_K^T(x) \to \mathbb{R}_{\geq 0}$ on
continuations, usable as one admissibility criterion among the family
that chapter organizes. It is emphatically not offered as the
criterion. A trajectory can maintain a large distinction margin while
being otherwise objectionable: slow, costly, fragile to small further
perturbation. The whole point of building a family of criteria
rather than a single one is that none of these concerns should be
allowed to silently stand in for the others.
\end{remark}

\section*{The manuscript, formally}

Return to the opening example with the machinery now in hand. Let $R$
be the set of ``intact page'' states, and let $\mathcal{D}$ be the
distinction separating pages bearing the original passage from pages
bearing any other passage of matching length and plausible style. Both
restorers produce trajectories landing in $R$: both are repair
solutions in the sense of the first section of this chapter, and
nothing said there distinguishes them. Only $A_{\mathcal D}$, evaluated
along each restorer's trajectory, can tell them apart. And it does:
the first restorer's trajectory, ending at a page whose
future reachable set (its readability, its citability, its
consequences for anyone who later relies on it) may be identical to
that of countless other plausible completions, will in general show
$A_{\mathcal D}$ collapsing toward zero; the second restorer's
trajectory, ending at the one page whose future is answerable to a
witness outside the manuscript itself, need not. Whether it actually
does depends on the specific $\mathcal D$ and dynamics in play, and no
general theorem is claimed here. The point of working the example
through is only to show that the formal apparatus, applied carefully,
draws the distinction the opening paragraph insisted mattered, rather
than washing it out.

\chapter{Admissibility Functionals}

A hiking guide choosing a trail for a mixed group of walkers has, in
front of her, several trails and several things she cares about: how
far each one runs, how steep it gets, how exposed it is to weather that
could turn in the afternoon, how good the views are for the members of
the group who came mainly for the views. It would be easy, and it
would be a mistake, to average these into a single score and rank the
trails by it. A score high on average can hide a trail that is
disastrous on the one dimension a particular walker cannot tolerate,
and a score low on average can hide a trail that is exactly right for
a group that only cares about two of the four criteria. The guide who
is good at her job does not collapse the criteria; she keeps them
separate for as long as she can, and only forces a single choice at
the last possible moment, when someone actually has to start walking.

The temptation to collapse many criteria into one number is old and
well-documented across the corpus this book is part of, and it is worth
resisting formally, not just rhetorically. What follows builds
admissibility, the ranking of continuations against more than mere
existence, in the order that resists the collapse for as long as
possible: first comparisons that need no number at all, then structures
built from several numbers kept separate, then the lattice those
structures naturally form, and only at the very end, as a special and
optional case, a single score. A single number is not banned from this
account. It is demoted to the last resort it should always have been.

\section*{Dominance relations}

\begin{definition}[Admissibility relation]
A dominance relation on $\Gamma_K^T(x)$ is a preorder $\preceq_A$
(reflexive and transitive, not necessarily total): $\gamma_1 \preceq_A
\gamma_2$ means $\gamma_2$ is at least as admissible as $\gamma_1$
according to standard $A$.
\end{definition}

\begin{remark}
This is deliberately the weakest structure that deserves the name.
It does not require every pair of continuations to be comparable:
two continuations can simply be incommensurable under $A$, neither
dominating the other. And it does not require a numerical score at
any point. Distinction preservation, from the previous chapter, is
already an instance: $\gamma_1 \preceq_{\mathcal D} \gamma_2$ iff
$A_{\mathcal D}(\gamma_1) \leq A_{\mathcal D}(\gamma_2)$ defines one
particular dominance relation, though nothing in this section requires
$A$ to be real-valued at all; a relation could compare continuations
directly, for instance by inclusion of the sets of distinctions each
one preserves, without ever reducing that comparison to a number.
\end{remark}

\section*{Pareto structures}

\begin{definition}[Pareto dominance]
Given criteria $A_1, \dots, A_n : \Gamma_K^T(x) \to \mathbb{R}$, define
$\gamma_1 \preceq_P \gamma_2$ iff $A_i(\gamma_1) \leq A_i(\gamma_2)$ for
every $i$.
\end{definition}

\begin{proposition}
$\preceq_P$ is a preorder on $\Gamma_K^T(x)$, and a genuine partial
order on the quotient by $\gamma_1 \sim \gamma_2 \iff A_i(\gamma_1) =
A_i(\gamma_2)\ \forall i$.
\end{proposition}
\begin{proof}
Reflexivity and transitivity hold coordinatewise, since each
$\leq$ on $\mathbb{R}$ is reflexive and transitive. Antisymmetry fails
on $\Gamma_K^T(x)$ itself whenever two distinct continuations achieve
identical criterion values, but holds by construction on the quotient
identifying such continuations.
\end{proof}

\begin{remark}
This already refuses the collapse the guide's instinct warned against:
$\gamma_1 \preceq_P \gamma_2$ requires $\gamma_2$ to be at least as good
on \emph{every} criterion, not merely on some weighted combination.
Trails that are better on distance but worse on exposure remain
genuinely incomparable under $\preceq_P$, exactly as they should.
\end{remark}

\section*{The admissibility lattice}

\begin{definition}
Write $\vec A(\gamma) := (A_1(\gamma), \dots, A_n(\gamma)) \in
\mathbb{R}^n$ for the criterion vector of $\gamma$.
\end{definition}

\begin{proposition}[$\mathbb{R}^n$ under $\preceq_P$ is a lattice]
For $\vec a, \vec b \in \mathbb{R}^n$, define
$\vec a \wedge \vec b := (\min(a_1,b_1), \dots, \min(a_n,b_n))$ and
$\vec a \vee \vec b := (\max(a_1,b_1), \dots, \max(a_n,b_n))$. Then
$\vec a \wedge \vec b$ is the greatest lower bound and
$\vec a \vee \vec b$ the least upper bound of $\{\vec a, \vec b\}$
under the coordinatewise order, so $(\mathbb{R}^n, \preceq_P)$ is a
lattice.
\end{proposition}
\begin{proof}
$\vec a \wedge \vec b \preceq_P \vec a$ and $\vec a \wedge \vec b
\preceq_P \vec b$ hold coordinatewise since $\min(a_i,b_i) \leq a_i$
and $\min(a_i,b_i) \leq b_i$. If $\vec c \preceq_P \vec a$ and
$\vec c \preceq_P \vec b$, then $c_i \leq a_i$ and $c_i \leq b_i$ for
every $i$, so $c_i \leq \min(a_i,b_i)$, giving
$\vec c \preceq_P \vec a \wedge \vec b$; hence $\vec a \wedge \vec b$
is the greatest lower bound. The argument for $\vee$ as least upper
bound is symmetric.
\end{proof}

\begin{remark}[The lattice structure does not automatically transfer to $\Gamma_K^T(x)$]
The proposition above is a fact about $\mathbb{R}^n$, not directly
about the space of continuations. Given two continuations $\gamma_1,
\gamma_2$, the vector $\vec A(\gamma_1) \wedge \vec A(\gamma_2)$ is a
well-defined point of $\mathbb{R}^n$, but there is no general reason to
expect some actual continuation $\gamma_3 \in \Gamma_K^T(x)$ achieving
exactly that criterion vector: the image
$\vec A(\Gamma_K^T(x)) \subseteq \mathbb{R}^n$ need not be closed under
$\wedge$ and $\vee$ even though $\mathbb{R}^n$ itself always is. When
the achievable region fails to be a sublattice of $\mathbb{R}^n$
(which can happen for entirely mundane reasons, such as the achievable
criterion vectors tracing out a non-convex curve rather than filling a
region), the lattice structure remains available as a way of
comparing criterion vectors abstractly, but stops being a promise that
some continuation realizes the meet or join of any two given ones.
\end{remark}

\section*{Scalar functionals as the last resort}

\begin{proposition}[Scalar functionals are total preorders]
Given $F : \Gamma_K^T(x) \to \mathbb{R}$, define $\gamma_1 \preceq_F
\gamma_2 \iff F(\gamma_1) \leq F(\gamma_2)$. Then $\preceq_F$ is a
total preorder: reflexive, transitive, and total (any two
continuations are comparable), with antisymmetry only up to
$F$-equivalence.
\end{proposition}
\begin{proof}
Reflexivity and transitivity are inherited from $\leq$ on $\mathbb{R}$.
Totality holds because any two real numbers $F(\gamma_1)$,
$F(\gamma_2)$ are comparable. Antisymmetry can fail on
$\Gamma_K^T(x)$ whenever distinct continuations achieve the same value
of $F$.
\end{proof}

\begin{remark}
A scalar functional is exactly a Pareto structure with $n=1$, or a
dominance relation that happens to be total. Nothing is wrong with
using one; Chapter 4's guaranteed value $V^\sharp$ is a scalar
functional and it did real work. What is being resisted here is
treating a scalar functional as the default starting point rather than
what it actually is: the special case obtained by discarding every
distinction between criteria before the comparison is made, appropriate
exactly when that discarding is a decision the modeler is willing to
own, not a shortcut taken to avoid building the Pareto structure first.
\end{remark}

\section*{The threshold of deferral}

Fix $K = \operatorname{Viab}_F(K)$ (the constraint set is already the
viability kernel, so its own topological boundary is the relevant
object), and suppose $K = \{g \leq 0\}$ for smooth $g$ with $\nabla g
\neq 0$ on $\partial K$, as in Chapter 3. Recall
$\delta(x) := \operatorname{dist}(x, X \setminus \operatorname{Viab}_F(K))$,
ambient distance to the non-viable region.

\begin{definition}[Tangential margin]
For $x \in K$ with $g(x) = 0$ (a boundary point),
\[
m(x) := -\min_{v \in F(x)} \nabla g(x) \cdot v = \max_{v \in F(x)} \langle -\nabla g(x), v \rangle,
\]
the depth to which the best available direction in $F(x)$ points into
the interior of $K$.
\end{definition}

By Chapter 3's smooth-boundary proposition, $x \in \partial K$ is
viable exactly when $m(x) \geq 0$; $m(x) = 0$ is the borderline case, a
direction in $F(x)$ exactly tangent to $\partial K$ and none pointing
further in.

\begin{proposition}[$\delta$ and $m$ are independent in general]
There exist systems in which $m(x)$ stays bounded away from $0$ as
$\delta(x) \to 0$, and systems in which $m(x) \to 0$ while $\delta(x)$
stays bounded away from $0$.
\end{proposition}
\begin{proof}
For the first: let $F(x) \equiv \{v_0\}$ for a fixed vector $v_0$ with
$\langle -\nabla g(x), v_0 \rangle \geq c > 0$ for all $x \in \partial K$
(a uniformly inward-pointing constant dynamics, possible whenever
$\partial K$ is, for instance, a hyperplane or has bounded curvature
relative to $v_0$). Then $m(x) \geq c > 0$ for every boundary point,
regardless of how close $x$ is taken to any other reference point, so
$m$ does not degenerate as $\delta(x) \to 0$ along $\partial K$ itself.

For the second: let $x_0$ be an interior point ($g(x_0) < 0$, so
$\delta(x_0) > 0$ strictly), and let $F(x_0) = \{0\}$ (the dynamics
vanishes at $x_0$ specifically, consistent with the standing
hypotheses if $F$ is otherwise well-behaved nearby). Then
$m(x_0)$ is not even defined by the boundary formula, but the
underlying tangential condition $F(x_0) \cap T_K(x_0)$ degenerates to
checking whether $0 \in T_{K}(x_0)$, which holds trivially since
$x_0$ is interior ($T_K(x_0) = X$ there); the relevant quantity is
instead the analogous margin against nearby non-viable directions
introduced by an adversarial perturbation, which can be made to vanish
at $x_0$ by construction while $x_0$ remains strictly interior to $K$,
i.e.\ $\delta(x_0)$ bounded away from $0$.
\end{proof}

\begin{remark}
This makes the appendix's caution fully rigorous rather than merely
plausible: ambient proximity to the boundary and degeneration of the
tangential margin are simply different facts about a system, and
neither implies the other in general. The heuristic status of $\delta$
as a stand-alone proxy is not a gap in this book's analysis; it is a
correct description of how these two quantities actually relate.
\end{remark}

\begin{proposition}[Local Lipschitz control near a critical boundary point]
Suppose $F$ is Hausdorff--Lipschitz in $x$ with constant $L$
($d_H(F(x),F(y)) \leq L|x-y|$), $\nabla g$ is Lipschitz with constant
$L_g$, and both $\sup_{v \in F(x)} \|v\|$ and $\|\nabla g(x)\|$ are
bounded by $C_F, C_g$ respectively on a neighborhood $U$ of
$\partial K$. Then $m$ is Lipschitz on $U \cap \partial K$ with constant
$C := C_F L_g + C_g L$. Consequently, if $x_0 \in \partial K$ is a
critical boundary point with $m(x_0) = 0$, then
\[
m(x) \leq C\,|x - x_0| \qquad \text{for all } x \in U \cap \partial K.
\]
\end{proposition}
\begin{proof}
Write $m(x) = h_{F(x)}(-\nabla g(x))$, where
$h_A(p) := \max_{a \in A} \langle p, a \rangle$ is the support function
of $A$. The support function satisfies the standard estimate
\[
|h_A(p) - h_B(q)| \leq \Big(\sup_{a \in A}\|a\|\Big)|p-q| + \|q\|\, d_H(A,B)
\]
for compact $A,B$ and $p,q \in X$ (triangle inequality via
$|h_A(p)-h_B(q)| \leq |h_A(p)-h_A(q)| + |h_A(q)-h_B(q)|$, bounding the
first term by the Cauchy--Schwarz-based Lipschitz estimate for support
functions in their argument, and the second by the standard bound
relating support functions to Hausdorff distance). Applying this with
$A = F(x)$, $p = -\nabla g(x)$, $B = F(y)$, $q = -\nabla g(y)$:
\[
|m(x) - m(y)| \leq C_F \|\nabla g(x) - \nabla g(y)\| + C_g\, d_H(F(x),F(y))
\leq C_F L_g |x-y| + C_g L |x-y| = C|x-y|.
\]
Taking $y = x_0$ and using $m(x_0) = 0$ gives $m(x) \leq C|x-x_0|$
directly.
\end{proof}

\begin{remark}
This is the positive result the appendix asked for, and it is worth
being precise about its actual reach. It does not say the tangential
margin degenerates near \emph{every} boundary point as ambient distance
shrinks: the first half of the independence proposition above rules
that out in general. It says that near a boundary point where the
margin is \emph{already zero}, Lipschitz continuity of the
dynamics and of the boundary's normal direction forces the margin to
shrink continuously as the point is approached, at a rate controlled
by $L$ and $L_g$, constants with a direct reading as, respectively,
how fast the admissible velocity set can change and how fast the
boundary's tangent plane can rotate, which is the
transversality-style control the appendix asked for. Whether a given
critical point exists at all, and where, remains a separate question
this proposition does not answer; what it answers is how the margin
must behave once such a point is known to exist. The general,
unconditional claim about $\delta$ as a stand-alone proxy remains
false, and should continue to be treated as false; what has changed is
that the conditional claim now has a proof rather than a citation to a
plausible-sounding hypothesis.
\end{remark}

\section*{Closing the family}

Nothing in this chapter selects among dominance relations, Pareto
structures, or scalar functionals as the one a given application
should use; that selection is a modeling decision made once, in the
open, for reasons stated outside the mathematics. What this chapter
has supplied is the machinery for making that decision honestly:
a vocabulary in which criteria can be kept separate for as long as
possible, a lattice that organizes them without silently ranking them,
and a precise account of when a single boundary quantity like $\delta$
can and cannot be trusted to stand in for the more detailed structure
underneath it. The next chapter extends this vocabulary from a single
system's continuations to several systems whose continuations
interact, where the question of whose criteria count, and how, cannot
be deferred any further.

\chapter{Preference Fields and Coordination Geometry}

Two ships can each be perfectly capable of crossing a narrow channel in
safety, and still not both be capable of crossing it at the same time.
Each captain, alone, faces an easy problem: plenty of room, plenty of
water, nothing in the way. Put both ships in the channel together and
the room does not change, but what each captain is entitled to do with
it does. Coordination failure of this kind is not a failure of either
captain's seamanship. It is a mismatch between what each of them can
achieve separately and what the channel can support when their claims
on it are added together.

Everything built so far in this book has quietly assumed a single
system asking what it can still become. Most of the situations worth
caring about involve more than one such system sharing a future:
departments in an organization, species in a habitat, players in a
game, coalitions in a negotiation. What changes, formally, when a
second decision-maker enters isn't the dynamics. The channel is
still the channel. What changes is the bookkeeping: each party now has its own
standard for what counts as an acceptable outcome, and those standards
must be checked not only against what the world allows but against
each other.

\section*{Correcting the minimal model}

Two different models could reasonably be meant by ``multiple agents
sharing a continuation space,'' and they are not interchangeable.
In one, each agent's control enters the state dynamics directly, so
that different agents can steer the system toward different reachable
states. In the other, the state dynamics are shared and
control-independent, driven only by an external perturbation, and what
differs between agents is only how they are paid along whatever
trajectory the shared dynamics produce. The appendix to this book
proposed a version closer to the first kind. Aubin's own construction
is the second kind. The difference is not cosmetic, and it is worth
having both models named before either is developed further.

The appendix to this book proposed, as a first attempt at a multi-agent
extension, a model in which each agent $i$ has its own constraint set
$K_i$ but all agents share one state and one dynamics $F$, with
$\Gamma_i(x) := \Gamma_{K_i}(x)$ and joint continuation given by
$\Gamma_{\mathrm{joint}}(x) := \bigcap_i \Gamma_i(x)$. That model is
usable, and the proposition proved from it, that individual viability
does not imply joint viability, is correct as far as it goes. But it
is not, in fact, the model Aubin's own
treatment of fuzzy dynamical cooperative games uses, and the difference
matters enough to correct before going further.

In Aubin's construction, the coalition's state dynamics depend only on
a shared perturbation, not on any agent's control at all:
$x'(t) = f(x(t), v(t))$, with $v(t) \in Q(x(t))$ the only source of
variation in how the state itself moves. Each agent's allotment
$p_i(t)$ enters only through a separate payoff equation attached to
that agent,
\[
y_i'(t) = \langle p_i(t), f(x(t),v(t)) \rangle - y_i(t)\, m_i(x(t), p_i(t), v(t)),
\]
governing how agent $i$'s own payoff accumulates along the shared
trajectory. There is, in this construction, only one continuation
space for the state itself: every agent shares exactly the same
$\Gamma_K(x)$, since no agent's choice of allotment changes which
states are reachable. What can fail to be shared is not the space of
possible trajectories, but whether every agent's payoff requirement can
be met simultaneously along one and the same trajectory, using
allotments that are themselves mutually constrained.

\begin{remark}
The earlier proposal is not wrong so much as answering a different
question: it describes a coherent and possibly useful generalization,
but a generalization \emph{beyond} Aubin's model, not a faithful
reconstruction of it. What follows reconstructs the model Aubin
actually built, since that is the one needed to use the fuzzy-games
paper directly, as promised.
\end{remark}

\section*{Joint admissibility under a shared resource}

\begin{definition}[Shared-state multi-agent continuation]
Fix $\Gamma_K(x)$ as in Chapter 1, shared by every agent
$i = 1, \dots, n$. Each agent has an allotment set
$P_i(x) \subseteq \mathbb{R}_+^k$ and a requirement functional
$b_i : \mathbb{R}_+ \times X \to \mathbb{R}_+$. Given
$\gamma \in \Gamma_K^T(x)$ and an allotment path
$p_i(\cdot) : [0,T] \to P_i(\gamma(\cdot))$, agent $i$'s payoff along
$\gamma$ is $y_i(t) := \langle p_i(t), \gamma(t) \rangle$, and agent
$i$'s admissibility criterion is
\[
A_i(\gamma, p_i(\cdot)) := \inf_{t \in [0,T]} \big[ y_i(t) - b_i(T-t, \gamma(t)) \big] \geq 0.
\]
\end{definition}

This is exactly the shape of an admissibility functional from the
previous chapter, now indexed by the agent's own choice of allotment
path as well as by the trajectory itself.

\begin{definition}[Joint feasibility under a shared resource]
Given a total available allotment $p(t) \in P(\gamma(t))$ (the grand
coalition's allotment, in Aubin's terms), a profile
$(p_1(\cdot), \dots, p_n(\cdot))$ is jointly feasible along $\gamma$ if
$\sum_i p_i(t) \preceq p(t)$ for all $t$ (coordinatewise, if allotments
are vector-valued), and $\gamma$ is jointly admissible if there exists
a jointly feasible profile with $A_i(\gamma, p_i(\cdot)) \geq 0$ for
every $i$ simultaneously.
\end{definition}

\begin{proposition}[Individual admissibility does not imply joint admissibility]
There exist systems in which every agent's requirement is individually
satisfiable along $\gamma$ (for each $i$ separately, some allotment
path achieves $A_i(\gamma,p_i(\cdot)) \geq 0$), while no jointly
feasible profile satisfies every agent's requirement at once.
\end{proposition}
\begin{proof}
Let $X = \mathbb{R}$, two agents, $\gamma(T) = 10$ fixed and
deterministic (no perturbation, for simplicity), $P_1 = P_2 = [0,1]$
scalar allotments, $b_1 = b_2 \equiv 0$ except at $t = T$ where each
agent requires $y_i(T) = p_i \cdot \gamma(T) \geq 6$. Individually:
agent 1 alone can take $p_1 = 0.6$, giving $y_1(T) = 6 \geq 6$; agent 2
alone can take $p_2 = 0.6$, giving $y_2(T) = 6 \geq 6$. Both
requirements are individually satisfiable. But the shared resource is
the grand coalition's allotment $p(T) = 1$ (the total available share
of $\gamma(T)$), so joint feasibility requires $p_1 + p_2 \leq 1$.
Satisfying both requirements simultaneously would require
$p_1 \geq 0.6$ and $p_2 \geq 0.6$, giving $p_1 + p_2 \geq 1.2 > 1$,
violating the resource constraint. Hence no jointly feasible profile
satisfies both agents' requirements.
\end{proof}

\begin{remark}
This is the coordination failure the two ships illustrated, made exact
and tied directly to Aubin's own economic content: it is a version of
the classical observation that a set of payoff demands can be
individually rational (each achievable in isolation) while exceeding
what a fixed total has to give, which is exactly the condition the
core of a cooperative game is built to characterize away
\cite{aubin2001}. The channel is not the scarce resource in this
telling; the grand coalition's total allotment is.
\end{remark}

\section*{Preference fields}

Admissibility, in the sense of the previous chapter, answers what is
allowed: which continuations satisfy the criteria imposed on them,
whether by a single system's own standards or, as above, by the
combined demands of several agents sharing a resource. It says
nothing about which of the allowed continuations any particular agent
would actually choose, given that more than one of them may qualify.

\begin{definition}[Preference field]
A preference field for agent $i$ is a map
$P_i : \Gamma_K^T(x) \to \mathbb{R}$ (or, more generally, a total
preorder $\preceq_{P_i}$), understood to be applied only to
continuations already admissible for agent $i$: $P_i$ ranks within
$\Gamma_i^\star(x) := \{\gamma : A_i(\gamma, p_i(\cdot)) \geq 0 \text{
for some feasible } p_i(\cdot)\}$, not across the whole of
$\Gamma_K^T(x)$.
\end{definition}

\begin{proposition}[Preference selects within admissibility]
If $\Gamma_i^\star(x)$ is nonempty and compact and $P_i$ is continuous,
then $\gamma_i^\ast := \operatorname{argmax}_{\gamma \in \Gamma_i^\star(x)} P_i(\gamma)$
exists and is agent $i$'s most preferred admissible continuation.
\end{proposition}
\begin{proof}
A continuous real-valued function attains its maximum on a nonempty
compact set (Weierstrass extreme value theorem), so
$\operatorname{argmax}_{\gamma \in \Gamma_i^\star(x)} P_i(\gamma)$ is
nonempty.
\end{proof}

\begin{remark}
This is the same $\Gamma \to \Gamma^\star$ pattern this book has used
since Chapter 4, applied once more, with one further layer: first the
menu $\Gamma$ is narrowed to the admissible subset $\Gamma^\star$ by
criteria that need not be agent-specific, and only then is a
genuinely agent-specific preference applied to pick out a single most-
wanted continuation from what remains admissible. Keeping the two
layers separate is the entire point: two agents facing identical
admissibility constraints can have entirely different preference
fields, and nothing about the admissibility layer needs to know or
care which agent is asking.
\end{remark}

\section*{Where the resource comes from}

Nothing in this chapter has said where the grand coalition's total
allotment $p(t)$ itself comes from, or how large it is allowed to be.
That is a further modeling choice, external to the mathematics
built here, just as the choice of admissibility criteria in the
previous chapter was external to it. What this chapter has supplied is
the machinery for stating, once that choice is made, exactly when a
group of agents with individually reasonable demands can and cannot be
jointly satisfied, and for keeping what each agent wants separate from
what any of them, or all of them together, are actually allowed to
have. The final chapter turns from this machinery to a handful of
worked sketches, in biology, in software, in institutions, showing
what the whole apparatus built across this book looks like applied to
systems that were never differential inclusions to begin with, only
things that had to keep going.

\chapter{Small Machines}

A physicist who has just written down a new differential equation
reaches for a pendulum before reaching for a galaxy. Not because
pendulums matter more than galaxies, but because a pendulum can be
built on a desk, swung by hand, and checked against the equation
within the hour, while a galaxy takes decades and a great deal of
faith to check against anything. The theory that survives contact with
the pendulum earns the right to be pointed at harder problems. The
theory that fails there was never going to survive the galaxy either;
it would only have taken longer to find out.

This book has, so far, mostly reached for galaxies: infinite-horizon
trajectory spaces, measures that turn out not to exist, boundaries
whose curvature is deliberately left undefined. All of that was
necessary, and none of it was checkable by hand. This chapter reaches
for the pendulum. It builds the smallest systems that still have
something to say: a circuit with two inputs and two outputs, a
counter with four states, a handful of bits standing in for a game
that has not been written yet, and asks the same questions of them
that the rest of this book has asked of everything else: what can this
system still become, how much room does it have, and does it remember
enough to know the difference between the futures that matter and the
ones that don't. Every answer in this chapter can be checked by hand,
which is exactly the point.

\section*{Combinational logic: continuation without a choice}

\begin{definition}[Half adder]
A half adder is the function $\varphi : \{0,1\}^2 \to \{0,1\}^2$,
$\varphi(a,b) = (a \oplus b,\ a \wedge b)$, the first output (sum) and
second output (carry) of adding two one-bit numbers.
\end{definition}

\begin{center}
\begin{tabular}{|c|c||c|c|}
\hline
$a$ & $b$ & sum $= a \oplus b$ & carry $= a \wedge b$ \\
\hline
0 & 0 & 0 & 0 \\
0 & 1 & 1 & 0 \\
1 & 0 & 1 & 0 \\
1 & 1 & 0 & 1 \\
\hline
\end{tabular}
\end{center}

Read as a continuation problem, $X = \{0,1\}^2$ (the input register),
and $\varphi$ is a one-step reachability map in exactly the sense of
Kauffman's adjacent possible from Chapter 2:
$R_1(a,b) := \{\varphi(a,b)\}$.

\begin{proposition}
For every $(a,b) \in \{0,1\}^2$, $R_1(a,b)$ is a singleton, and hence
$\operatorname{Vol}_1(a,b) = 1$ under counting measure.
\end{proposition}
\begin{proof}
Immediate from the truth table: $\varphi$ is a function, not a
relation, so it assigns exactly one output pair to each input pair.
\end{proof}

\begin{remark}
This is the degenerate case of every construction in this book, and it
is worth pausing on exactly how degenerate. There is no viability
question to ask (every state trivially has a successor, so
$\operatorname{Viab}_F(X) = X$ with no boundary at all), no volume
question worth asking (every reachable set has size one, so
$\operatorname{Vol}_1$ carries no information), and admissibility
collapses to the strictest possible case: a single output is either
exactly correct or it is a fault, with nothing in between and no
lattice of criteria to build. A combinational circuit is, formally, a
continuation problem with the continuation stripped out: one step,
one outcome, no room to move. The elaborate machinery built across
this book correctly predicts its own irrelevance here, which is a
better sign than it might first appear: a theory that cannot recognize
a trivial case is not to be trusted with a hard one.
\end{remark}

\begin{remark}
This is also the cleanest possible instance of the Markovian case from
Chapter 5. A combinational circuit has no history-dependence to
resolve because it has no history at all: $\varphi(a,b)$ depends only
on the current input register, never on how that register came to hold
the values it holds. The sufficient statistic is the empty statistic,
$\varphi(H_t) := \varnothing$, trivially autonomous, trivially
sufficient. Sequential circuits, taken up next, are what happens when
that stops being true.
\end{remark}

\section*{Sequential circuits: continuation with a memory}

\begin{definition}[Enabled counter]
Let $X = \{0,1,2,3\}$ (a 2-bit register), and let the one-step map
depend on an enable input $e \in \{0,1\}$:
\[
\Phi(x, e) = \begin{cases} \{x\} & e = 0 \\ \{(x+1) \bmod 4\} & e = 1. \end{cases}
\]
\end{definition}

Unlike the half adder, this machine has a genuine reachable-set
question to ask, because the enable input is not fixed in advance:
treating $e$ as a control available at every step gives
$F(x) := \{x, (x+1) \bmod 4\}$, a genuine two-element reachable set at
every state.

\begin{proposition}
Starting from $x_0 = 0$: $\operatorname{Vol}_0(0) = 1$,
$\operatorname{Vol}_1(0) = 2$, $\operatorname{Vol}_2(0) = 3$,
$\operatorname{Vol}_3(0) = 4$, and $\operatorname{Vol}_T(0) = 4$ for
all $T \geq 3$.
\end{proposition}
\begin{proof}
$R_0(0) = \{0\}$. $R_1(0) = F(0) = \{0,1\}$. $R_2(0) = F(0) \cup F(1) =
\{0,1\} \cup \{1,2\} = \{0,1,2\}$. $R_3(0) = \{0,1,2\} \cup F(2) =
\{0,1,2\} \cup \{2,3\} = \{0,1,2,3\}$, all four states of $X$. Since
$X$ has only four elements, $R_T(0) \subseteq X$ can grow no further,
and since every state has itself in its own reachable set (via
$e = 0$), $R_T(0)$ never shrinks once it reaches $X$; hence
$R_T(0) = X$ for all $T \geq 3$.
\end{proof}

\begin{remark}
This is volume behaving the way intuition expects: monotone
nondecreasing, saturating once the whole space is reached, in
contrast to the folding, non-monotone example of Chapter 6. The
difference is structural: this counter's transitions include a
self-loop at every state ($e=0$ always available), which rules out the
kind of collapse that produced non-monotonicity there. A counter
\emph{without} a stay-put option ($\Phi(x) = \{(x+1) \bmod 4\}$ only,
no choice at all) reproduces Chapter 6's oscillating example exactly:
$R_T(0)$ would be a single point cycling through $0,1,2,3,0,1,2,\dots$,
with $\operatorname{Vol}_T(0) = 1$ forever, never growing, a half
adder's worth of freedom stretched over four states instead of one
step.
\end{remark}

\begin{remark}[The counter as its own sufficient statistic]
The current count is a sufficient statistic for this system's
continuation operator in the exact sense of Chapter 5: $\Gamma(x_t \mid
H_t)$ depends on $H_t$ only through $x_t$ itself, since $\Phi$ takes
only the current count and the current enable input, never anything
earlier. This is close to the smallest nontrivial illustration
available of a system for which the trivial embedding of Chapter 5 and
the genuine finite reduction coincide: the sufficient statistic and the
state are, here, the same two bits.
\end{remark}

\section*{Toward eight bits: a state machine with something to lose}

A two-bit counter has no distinctions worth preserving, because it has
nothing riding on its state beyond the count itself. A byte used to
represent something (a game entity's position and status, an
oscillator's phase and amplitude) can fail in a way a bare counter
cannot: two logically different situations can come to share the same
bit pattern, at which point the system has no way left to tell them
apart. This is not a hypothetical concern. Constrained-memory systems,
of the kind eight-bit hardware forced on a generation of
programmers, have a well-documented bug class arising from
this failure: states that were supposed to remain distinguishable
quietly colliding because there were not enough bits left to keep them
apart, producing behavior that looks inexplicable until the bit layout
is inspected directly. This is distinction collapse, in the exact
sense of Chapter 7, occurring in hardware small enough to hold in one
hand.

\begin{example}[A four-state traffic light with a memory bit]
Let $X = \{0,1,2,3\} \times \{0,1\}$: two bits for light phase (red,
red-yellow, green, yellow, cycling in that order) and one bit recording
whether a pedestrian has requested crossing. Transitions advance the
phase deterministically, $\Phi((x,r), \cdot) = ((x+1) \bmod 4,\ r')$,
where the request bit $r'$ is cleared ($r' = 0$) if the new phase is
red (phase $0$, the point at which a pedestrian may cross) and
otherwise carried forward ($r' = r$) unless a new request arrives.
\end{example}

\begin{definition}
Let $\mathcal{D}$ be the distinction separating states with a pending
request ($r=1$) from states without one ($r=0$), restricted to a fixed
phase.
\end{definition}

\begin{proposition}
$\mathcal{D}$ is preserved by $\Phi$ at every phase except phase $0$,
where it collapses: $A_{\mathcal D}^1((0,1)) = 0$, since
$\Phi((0,1),\cdot) = (1,0) = \Phi((0,0),\cdot)$.
\end{proposition}
\begin{proof}
At phase $0$, the transition rule clears the request bit regardless of
its incoming value, so $(0,1)$ and $(0,0)$ map to the identical state
$(1,0)$: their one-step reachable sets coincide, $R_1((0,0)) =
R_1((0,1)) = \{(1,0)\}$, giving $d_\Gamma^1((0,0),(0,1)) = 0$ and hence
$A_{\mathcal D}^1((0,1)) = 0$ by definition. At any other phase
$k \neq 0$, the request bit is carried forward unchanged, so
$\Phi((k,0),\cdot) = ((k+1)\bmod 4, 0) \neq ((k+1) \bmod 4, 1) =
\Phi((k,1),\cdot)$, giving distinct reachable sets and hence
$A_{\mathcal D}^1((k,1)) > 0$.
\end{proof}

\begin{remark}
This is exactly correct behavior for a traffic light: the request is
supposed to be consumed once crossing becomes possible. But the same formal signature, a distinction collapsing
at a designated point, is indistinguishable in form from a genuine bug:
a byte in a game engine that is supposed to remember which of two
enemies dealt the last hit to the player, but happens to get
overwritten by an unrelated update at a particular frame, collapses in
exactly the same way, for exactly the same reason, and the machinery
built in Chapter 7 does not know the difference between an intended
reset and an unintended one. Telling them apart is a question about
what the designer wanted preserved, not a question this chapter's
mathematics can answer; what the mathematics can do is locate, exactly
and by direct computation, every point at which a chosen distinction
does or does not survive, which is the diagnostic half of the problem
and the half that is easy to get wrong by eye once the state space
grows from six configurations to two hundred and fifty-six.
\end{remark}

\begin{remark}[Scaling to a full byte]
Nothing above depended on the state space being small other than
tractability of writing out the transition table by hand. A full
eight-bit register, $X = \{0,1\}^8$, with an update rule built the same
way (some bits advancing a phase or a position, others recording
flags whose job is to persist until a specific condition clears them)
is analyzed by the same questions, computed the same
way, just over $256$ states instead of $8$. Is the space of legal
states viable under the update rule: does the machine have a strategy
that never produces a state outside its intended range, a crash or
an out-of-bounds sprite in the vocabulary of a game engine? How much
room does a given state have to become something else before the level
or the phrase ends: continuation volume, now genuinely informative
rather than degenerate, since branching inputs make $\operatorname{Vol}_T$
vary meaningfully across states. And which flags survive which
transitions: distinction preservation, checked bit by bit rather than
by playtesting and hoping. None of this requires simulating every one
of the $256$ states by hand, any more than Chapter 6 required
simulating every trajectory of a continuous system by hand. It requires
only that the reachable-set and distinction-margin computations be
mechanical enough to hand to a computer once the update rule is
written down precisely, which a truth table always is, and a
differential inclusion isn't.
\end{remark}

\section*{What the small machines bought}

Every question this book has asked about infinite-dimensional
trajectory spaces has an exact, hand-checkable counterpart on a state
space small enough to write out in full: viability becomes a
reachability check on a finite graph, volume becomes counting, distance
becomes set difference, distinction margin becomes an equality test on
successor states. None of the open problems flagged earlier in this
book, the nonexistent trajectory-space measure, the unresolved status
of curvature, the case-by-case nature of sufficient statistics, are
open here, because finiteness resolves them for free: there is only
one sensible measure on a finite set, and it was never in question.
What the small machines cannot do is stand in for the continuous and
infinite-horizon cases this book spent most of its length on; a half
adder is a proof that the formalism is coherent on easy ground, not a
proof that it says anything true on hard ground. It is, in exactly the
sense the opening of this chapter intended, the pendulum, not the
galaxy. The conclusion that follows is about what has and has not been
earned across both.

\chapter{Closing Sketches, and What Was Earned}

A book that opens by proposing a central object owes its reader, at
the end, an honest answer to the only question that ever really
mattered: did the object hold up. Elegance is cheap, and this book has
occasionally indulged in it, so that's not really the question. The
real one is whether the thing proposed in the first chapter, the fan
of futures reachable from a state, turned out to be sturdy enough to
carry everything that got built on top of it, or whether, at some
point along the way, it had to be quietly propped up by something else
smuggled in from outside. The honest answer, worked out chapter by
chapter rather than asserted here, is: mostly, and not without help.
This closing chapter says what the help was, sketches three places
where the whole
apparatus is likely to go next, and then stops.

\section*{Three seeds}

None of what follows is a treatment. Each is a paragraph's worth of
translation, showing that the vocabulary built across this book has
something to say about a domain it was not built from, and each is
offered as the beginning of an independent monograph rather than a
compressed version of one. That's consistent with how Repair Theory, Fiscal
Reachability, and the Ecology of Distinctions each grew, in their own
time, out of a chapter that was originally meant to be brief.

\begin{remark}[Biological repair]
A regenerating organism is, in the vocabulary of this book, a system
whose viable configurations $K$ include not just its current anatomy
but every anatomy compatible with continued function, and whose repair
target $R$ is a specific one among them: the target morphology a
damaged tissue is, remarkably, often still able to reach. Levin's
bioelectric memory, discussed in Chapters 2 and 5 as the clearest
existing example of genuine history-dependence, is naturally read as
the physical carrier of a sufficient statistic, not the full cellular
history of the tissue but some lower-dimensional bioelectric pattern
sufficient to determine what the tissue can still become. Whether that
pattern admits a sufficient statistic in the sense Chapter 5 made
precise, and in particular whether it can be a genuinely
finite-dimensional one (a question that chapter's impossibility result
left explicitly open rather than settled), is an empirical question
this book does not answer and a formal one a future volume could.
\end{remark}

\begin{remark}[Software systems]
A codebase's continuation space is the set of behaviors reachable from
its current source by some sequence of edits that keeps the tests
passing. $K$ is, in an ordinary sense, what a test suite is
for. A refactor is a repair problem with $R$ set to ``same observable
behavior, different internal structure.'' Technical debt is a
diminishing continuation volume: not fewer lines of code, but fewer
future edits that remain safe to make without breaking something a
test does not check for. And a regression is a distinction
collapse in the sense of Chapter 7: two states of the system, one
correct and one subtly incorrect, that a test suite was supposed to keep apart,
quietly merging under an edit that the suite was not built to notice.
Diagnostic coordinates, elsewhere in this corpus, are close to
the margin $A_{\mathcal D}$ of Chapter 7 applied to this domain: how
much information remains available to tell a correct state from a
broken one before the gap disappears entirely.
\end{remark}

\begin{remark}[Institutions]
An institution's viable configurations are the arrangements under
which it continues to perform the functions it exists to perform;
its capture basin, when those functions are already failing, is the
set of reforms from which recovery remains possible before it is not.
Chapter 9's coordination geometry is the most directly relevant part
of this book to institutions specifically, since institutions are
almost never single agents: a reform that is individually reasonable
to every stakeholder can still be jointly infeasible under a shared
resource, budget, legitimacy, public attention, in just the way
Chapter 9's two agents could each individually afford a payoff neither
could jointly receive. What this book does not supply, and what an
institutional volume would have to, is any account of where the
distinctions worth preserving in a civic system come from in the first
place. That's a harder and more contested question than anything Chapter 7
had to answer for a manuscript or a traffic light.
\end{remark}

\section*{What was established}

Stated plainly, without the hedges that earned it: this book proved
Aubin's Viability Theorem in full, both directions, under standing
hypotheses stated once and used throughout. It proved that the
guaranteed value already inside viability theory is a genuine
selection, not merely an existence witness, with a fully worked
example showing the selection can be strict and consequential rather
than vacuous. It resolved, to the extent either question is honestly
resolvable in general, when history-dependence can be reduced to a
larger but still Markovian state, and proved by explicit construction
that the reduction is not always available. It built a workable
continuation volume on finite horizons after proving, rather than
merely asserting, that the naive infinite-horizon version is
identically zero under any Wiener-type measure. It gave the word
distinction a formal meaning precise enough to prove exactly when a
distinction collapses under continuation, tied directly to the
distance already built rather than introduced as a separate notion. It
organized admissibility in an order that resists collapsing many
criteria into one for as long as possible, and proved a genuine, if
strictly local and conditional, connection between ambient boundary
distance and the degeneration of tangential margin, catching and
correcting its own earlier overreach on that same question along the
way. It caught and corrected a second overreach, in its own proposed
multi-agent model, once closer contact with Aubin's actual construction
showed the correction was needed. And it checked all of this, at the
smallest possible scale, against machines built from truth tables
rather than differential inclusions, where every claim could be
verified by hand rather than taken on faith.

\section*{What was not}

The naive continuation volume was not repaired, only replaced; nobody
in these pages built a genuine measure on an infinite-dimensional
trajectory space, and the finite-horizon workaround, while sound, is
an approximation whose relationship to whatever infinite-horizon
object it approximates was characterized only through the nested
intersection of Chapter 3, not through any direct construction.
Curvature was never defined, on the explicit and repeated grounds that
inventing a definition to fit a word already in use elsewhere would
have been worse than leaving the word unclaimed. The choice among
Hausdorff, Gromov--Hausdorff, Wasserstein, and reachability-based
distances was never made, because no worked example in this book
needed to make it, and a choice made without need is a choice made
badly. The general question of when a finite sufficient statistic
exists was answered only in the negative, by exhibiting a system
with none, and the positive cases were handled one at a time, as the
appendix always warned they would have to be. And the candidate for a
genuine unifying theorem, functoriality of the $\Gamma \to
\mathcal{S}(\Gamma) \to \Gamma^\star$ hierarchy under morphisms of
constrained systems, was named at the outset of this book and never
attempted; it remains exactly where it was left, a real question with
no false progress claimed against it.

\section*{Was $\Gamma$ enough}

This was the question the first chapter promised to test, and the
honest answer is neither yes nor no. $\Gamma(x)$, the bare set of
futures reachable from a state, was enough to organize existence,
enough to host a selection criterion once one was supplied from
outside, enough to carry a volume and a distance once the right
finite-horizon restriction was found, and enough to make distinction
collapse a checkable fact rather than an intuition. It was not enough,
on its own, to say which futures a history should be allowed to
depend on, which distinctions were worth drawing in the first place,
or how much of a shared resource any one among several agents was
entitled to. Each of those needed something brought in from outside
$\Gamma$ itself: a sufficient statistic, a partition, a resource
constraint. The actual content of this book, chapter by chapter,
was less the elegance of $\Gamma$ than the precision with which each of
those outside additions could be named, bounded, and in two cases
caught in error and fixed. A central object that needs help at almost
every turn, and is honest about exactly what help it needed and
exactly where, is not a failure of the object. It is what a central
object is for.

\appendix

\chapter*{Appendices}
\addcontentsline{toc}{chapter}{Appendices}

Several propositions in the main text cited a result as ``standard''
and moved on, in the interest of keeping each chapter's argument
readable rather than self-contained. That is a defensible choice in
context, but it leaves debts, and these three appendices pay them:
the support-function estimate used in Chapter 8, the measurable
selection step used in Chapter 4 and implicit throughout Chapter 3,
and the Gr\"onwall-type a priori bounds compressed into a phrase in
two separate proofs in Chapter 3. Each appendix is self-contained and
can be read independently of the others.

\chapter{Support Functions and Lipschitz Estimates}

\begin{definition}
For a nonempty compact $A \subseteq \mathbb{R}^n$, the support function
of $A$ is $h_A(p) := \max_{a \in A} \langle p, a \rangle$, for
$p \in \mathbb{R}^n$. (The maximum is attained since $A$ is compact and
$a \mapsto \langle p,a \rangle$ is continuous.)
\end{definition}

\begin{lemma}[Lipschitz continuity in the argument]
\label{lem:hlip}
$|h_A(p) - h_A(q)| \leq R_A\, |p-q|$, where $R_A := \max_{a \in A} \|a\|$.
\end{lemma}
\begin{proof}
Let $a^\ast \in A$ achieve $h_A(p) = \langle p, a^\ast \rangle$. Then
\[
h_A(p) - h_A(q) = \langle p, a^\ast \rangle - h_A(q) \leq \langle p, a^\ast \rangle - \langle q, a^\ast \rangle = \langle p-q, a^\ast \rangle \leq \|p-q\|\,\|a^\ast\| \leq \|p-q\|\, R_A,
\]
using $h_A(q) \geq \langle q, a^\ast \rangle$ (since $a^\ast \in A$) for
the first inequality and Cauchy--Schwarz for the second. Exchanging the
roles of $p$ and $q$ gives $h_A(q) - h_A(p) \leq R_A\|p-q\|$ by the same
argument, so $|h_A(p)-h_A(q)| \leq R_A\|p-q\|$.
\end{proof}

\begin{lemma}[Bound via Hausdorff distance]
\label{lem:hhaus}
For nonempty compact $A, B \subseteq \mathbb{R}^n$ and $q \in \mathbb{R}^n$,
$|h_A(q) - h_B(q)| \leq \|q\|\, d_H(A,B)$.
\end{lemma}
\begin{proof}
Let $a^\ast \in A$ achieve $h_A(q) = \langle q, a^\ast \rangle$. Since
$d_H(A,B) = \max\{\sup_{a \in A} \operatorname{dist}(a,B),\ \sup_{b \in B} \operatorname{dist}(b,A)\}$,
there is $b \in B$ with $\|a^\ast - b\| \leq d_H(A,B)$ (attained, by
compactness of $B$). Then
\[
h_B(q) \geq \langle q, b \rangle = \langle q, a^\ast \rangle + \langle q, b - a^\ast \rangle \geq h_A(q) - \|q\|\,\|b-a^\ast\| \geq h_A(q) - \|q\|\, d_H(A,B),
\]
so $h_A(q) - h_B(q) \leq \|q\|\,d_H(A,B)$. Exchanging the roles of $A$
and $B$ gives $h_B(q) - h_A(q) \leq \|q\|\,d_H(A,B)$ by the same
argument, so $|h_A(q)-h_B(q)| \leq \|q\|\,d_H(A,B)$.
\end{proof}

\begin{proposition}[The estimate used in Chapter 8]
For nonempty compact $A,B$ and $p,q \in \mathbb{R}^n$,
\[
|h_A(p) - h_B(q)| \leq R_A\,|p-q| + \|q\|\, d_H(A,B).
\]
\end{proposition}
\begin{proof}
By the triangle inequality, $|h_A(p)-h_B(q)| \leq |h_A(p)-h_A(q)| +
|h_A(q)-h_B(q)|$. Bound the first term by Lemma~\ref{lem:hlip} and the
second by Lemma~\ref{lem:hhaus}.
\end{proof}

\begin{remark}
This is exactly the estimate invoked without proof in Chapter 8's
Lipschitz-continuity proposition for the tangential margin $m(x) =
h_{F(x)}(-\nabla g(x))$: taking $A = F(x)$, $B = F(y)$,
$p = -\nabla g(x)$, $q = -\nabla g(y)$ recovers the bound used there
verbatim.
\end{remark}

\chapter{Measurable Selections}

Several proofs in this book pass silently between two descriptions of
the same trajectory: as a curve satisfying a differential inclusion,
$\dot\gamma(t) \in F(\gamma(t))$, and as a curve driven by an explicit,
individually chosen control or perturbation, $\dot\gamma(t) =
f(\gamma(t), v(t))$ for some measurable $v(t) \in Q(\gamma(t))$. Moving
from the first description to the second requires selecting, at each
time $t$, one particular element of the set
$\{v \in Q(\gamma(t)) : f(\gamma(t),v) = \dot\gamma(t)\}$, and doing so
in a way that is itself measurable in $t$: not merely nonempty at
each instant, which the inclusion already guarantees, but selectable
without discontinuous jumps that would break integrability.

\begin{theorem}[Kuratowski--Ryll-Nardzewski selection theorem, cited form]
Let $(T,\Sigma)$ be a measurable space and let $\Phi : T \rightrightarrows \mathbb{R}^n$
be a set-valued map with nonempty closed values such that
$\{t : \Phi(t) \cap U \neq \varnothing\} \in \Sigma$ for every open
$U \subseteq \mathbb{R}^n$ (weak measurability). Then $\Phi$ admits a
measurable selection: a measurable function $\sigma : T \to \mathbb{R}^n$
with $\sigma(t) \in \Phi(t)$ for every $t$.
\end{theorem}

\begin{proof}
This is the Kuratowski--Ryll-Nardzewski theorem, a foundational result
of descriptive set theory and set-valued analysis \cite{krn1965}; its
proof is a transfinite or countable successive-refinement construction
that lies outside the scope of this book and is not reproduced here.
\end{proof}

\begin{proposition}[Filippov-type selection for the applications in this book]
\label{prop:filippov}
Let $f : X \times V \to X$ be continuous, $Q : X \rightrightarrows V$
weakly measurable with nonempty closed values, and $F(x) := \{f(x,v) :
v \in Q(x)\}$. Let $\gamma$ be absolutely continuous with
$\dot\gamma(t) \in F(\gamma(t))$ for a.e.\ $t$. Then there exists a
measurable $v(\cdot)$ with $v(t) \in Q(\gamma(t))$ and
$\dot\gamma(t) = f(\gamma(t), v(t))$ for a.e.\ $t$.
\end{proposition}
\begin{proof}[Proof idea]
Define $\Phi(t) := \{v \in Q(\gamma(t)) : f(\gamma(t),v) = \dot\gamma(t)\}$,
nonempty for a.e.\ $t$ by hypothesis. Continuity of $f$ and
measurability of $\gamma$, $\dot\gamma$, and $Q(\gamma(\cdot))$ together
give weak measurability of $t \mapsto \Phi(t)$ (a standard consequence
of the Kuratowski--Ryll-Nardzewski framework, argued by showing the
graph of $\Phi$ is a measurable subset of $T \times V$ and invoking the
projection and selection theorems associated with it
\cite{aubinframkowska1990}). The Kuratowski--Ryll-Nardzewski theorem
then supplies a measurable selection $v(t) \in \Phi(t)$, which by
construction satisfies $f(\gamma(t),v(t)) = \dot\gamma(t)$ and
$v(t) \in Q(\gamma(t))$ for a.e.\ $t$.
\end{proof}

\begin{remark}
This is the result invoked, by name only, in Chapter 4's proof that
uncontrolled viability is the trivial case of the controlled framework,
and it underlies every passage in Chapters 3 and 6 that moves between
``a trajectory of the inclusion'' and ``a trajectory driven by some
admissible control or perturbation.'' The genuinely hard mathematics is
entirely inside the Kuratowski--Ryll-Nardzewski theorem itself; what
this appendix supplies is only the (still nontrivial, but comparatively
routine) argument that the specific set-valued maps arising in this
book satisfy that theorem's hypotheses.
\end{remark}

\chapter{Gr\"onwall-Type A Priori Bounds}

\begin{lemma}[Continuous Gr\"onwall inequality]
\label{lem:gronwall}
Let $w : [0,M] \to \mathbb{R}_{\geq 0}$ be continuous and satisfy
$w(t) \leq w(0) + c\int_0^t w(s)\, ds$ for a constant $c \geq 0$ and all
$t \in [0,M]$. Then $w(t) \leq w(0)\, e^{ct}$ for all $t \in [0,M]$.
\end{lemma}
\begin{proof}
Let $\Psi(t) := w(0) + c\int_0^t w(s)\,ds$, so $w(t) \leq \Psi(t)$ and
$\Psi'(t) = c\,w(t) \leq c\,\Psi(t)$ (using $w \leq \Psi$ and $c \geq 0$).
Then $\frac{d}{dt}\big(\Psi(t) e^{-ct}\big) = e^{-ct}(\Psi'(t) - c\Psi(t)) \leq 0$,
so $\Psi(t)e^{-ct}$ is nonincreasing, giving
$\Psi(t)e^{-ct} \leq \Psi(0) = w(0)$, i.e.\ $\Psi(t) \leq w(0)e^{ct}$.
Since $w(t) \leq \Psi(t)$, the result follows.
\end{proof}

\begin{proposition}[A priori bound under linear growth]
\label{prop:apriori}
Suppose $\gamma$ satisfies $\dot\gamma(t) \in F(\gamma(t))$ a.e.\ with
$\sup\{\|v\| : v \in F(x)\} \leq c(1+\|x\|)$ for all $x$ (hypothesis
(H3) of Chapter 3). Then for all $t \geq 0$,
\[
\|\gamma(t)\| \leq (1 + \|\gamma(0)\|)\, e^{ct} - 1.
\]
\end{proposition}
\begin{proof}
Let $w(t) := 1 + \|\gamma(t)\|$. Since
$\|\gamma(t)\| \leq \|\gamma(0)\| + \int_0^t \|\dot\gamma(s)\|\, ds \leq
\|\gamma(0)\| + \int_0^t c(1+\|\gamma(s)\|)\, ds = \|\gamma(0)\| + c\int_0^t w(s)\, ds$,
adding $1$ to both sides gives $w(t) \leq w(0) + c\int_0^t w(s)\, ds$.
By Lemma~\ref{lem:gronwall}, $w(t) \leq w(0)e^{ct} = (1+\|\gamma(0)\|)e^{ct}$,
so $\|\gamma(t)\| = w(t) - 1 \leq (1+\|\gamma(0)\|)e^{ct} - 1$.
\end{proof}

\begin{remark}
This is the uniform bound invoked as ``a discrete Gr\"onwall argument''
in the sufficiency proof of the Viability Theorem and again in the
proof that the viability kernel is a nested intersection of
finite-horizon sets, in both cases used to establish that a sequence
of admissible trajectories is uniformly bounded on any fixed interval
$[0,M]$, which is what licenses the Arzel\`a--Ascoli step in each
proof. The bound here depends only on $\|\gamma(0)\|$, $c$, and $t$,
uniformly across every trajectory satisfying the same growth
hypothesis, exactly the uniformity those proofs needed and asserted
without deriving.
\end{remark}

\begin{lemma}[Discrete Gr\"onwall inequality]
Let $a_0, a_1, a_2, \dots \geq 0$ satisfy $a_{k+1} \leq (1+hc)\,a_k + hc$
for constants $h, c > 0$. Then $a_k \leq (1+a_0)(1+hc)^k - 1$ for all
$k \geq 0$.
\end{lemma}
\begin{proof}
By induction. The case $k=0$ is immediate. Assume
$a_k \leq (1+a_0)(1+hc)^k - 1$. Then
\[
a_{k+1} \leq (1+hc)a_k + hc \leq (1+hc)\big[(1+a_0)(1+hc)^k - 1\big] + hc
= (1+a_0)(1+hc)^{k+1} - (1+hc) + hc,
\]
and $-(1+hc)+hc = -1$, giving $a_{k+1} \leq (1+a_0)(1+hc)^{k+1} - 1$.
\end{proof}

\begin{remark}
This is the discrete analogue used in the Euler polygon construction of
the same sufficiency proof, where the piecewise-linear approximate
trajectories $x_0, x_1, x_2, \dots$ satisfy exactly this kind of
recursive bound once the growth hypothesis is applied at each step, and
$(1+hc)^k \to e^{ckh}$ as $h \to 0$ with $kh$ fixed recovers the
continuous bound of Proposition~\ref{prop:apriori} in the limit, which
is what allows the discrete and continuous versions of the argument to
agree in the passage to the limit $h \to 0$.
\end{remark}

\chapter*{References}
\addcontentsline{toc}{chapter}{References}

\begin{thebibliography}{99}

\bibitem{aubin2001} J.-P.\ Aubin, \emph{Dynamic Core of Fuzzy Dynamical
Cooperative Games}, Annals of Dynamic Games, 2001.

\bibitem{aubin1991} J.-P.\ Aubin, \emph{Viability Theory}, Birkh\"auser,
1991.

\bibitem{aubincellina1984} J.-P.\ Aubin, A.\ Cellina, \emph{Differential
Inclusions}, Springer-Verlag, 1984.

\bibitem{aubinframkowska1990} J.-P.\ Aubin, H.\ Frankowska,
\emph{Set-Valued Analysis}, Birkh\"auser, 1990.

\bibitem{alexander1977} C.\ Alexander, S.\ Ishikawa, M.\ Silverstein,
\emph{A Pattern Language}, Oxford University Press, 1977.

\bibitem{bellman1957} R.\ Bellman, \emph{Dynamic Programming}, Princeton
University Press, 1957.

\bibitem{filippov1988} A.\ F.\ Filippov, \emph{Differential Equations
with Discontinuous Righthand Sides}, Kluwer Academic Publishers, 1988.

\bibitem{filippov1962} A.\ F.\ Filippov, ``On certain questions in the
theory of optimal control,'' SIAM Journal on Control, 1 (1962).

\bibitem{gronwall1919} T.\ H.\ Gr\"onwall, ``Note on the derivatives
with respect to a parameter of the solutions of a system of
differential equations,'' Annals of Mathematics, 20 (1919).

\bibitem{holland1995} J.\ Holland, \emph{Hidden Order: How Adaptation
Builds Complexity}, Addison-Wesley, 1995.

\bibitem{kauffman2000} S.\ Kauffman, \emph{Investigations}, Oxford
University Press, 2000.

\bibitem{krn1965} K.\ Kuratowski, C.\ Ryll-Nardzewski, ``A general
theorem on selectors,'' Bulletin de l'Acad\'emie Polonaise des
Sciences, 13 (1965).

\bibitem{levin2021} M.\ Levin, ``Bioelectric signaling: Reprogrammable
circuits underlying embryogenesis, regeneration, and cancer,'' Cell,
184(8), 2021.

\bibitem{lyapunov1892} A.\ M.\ Lyapunov, \emph{The General Problem of
the Stability of Motion}, 1892 (trans.\ Taylor \& Francis, 1992).

\bibitem{paleywienerzygmund1933} R.\ Paley, N.\ Wiener, A.\ Zygmund,
``Notes on random functions,'' Mathematische Zeitschrift, 37 (1933).

\bibitem{pearl2000} J.\ Pearl, \emph{Causality: Models, Reasoning, and
Inference}, Cambridge University Press, 2000.

\bibitem{poincare1890} H.\ Poincar\'e, \emph{Sur le probl\`eme des trois
corps et les \'equations de la dynamique}, Acta Mathematica, 13 (1890).

\bibitem{shreve2004} S.\ Shreve, \emph{Stochastic Calculus for Finance
II: Continuous-Time Models}, Springer, 2004.

\bibitem{utkin1992} V.\ Utkin, \emph{Sliding Modes in Control and
Optimization}, Springer-Verlag, 1992.

\end{thebibliography}

\end{document}
